%%
%% This is file `sample-manuscript.tex',
%% generated with the docstrip utility.
%%
%% The original source files were:
%%
%% samples.dtx  (with options: `all,proceedings,bibtex,manuscript')
%% 
%% IMPORTANT NOTICE:
%% 
%% For the copyright see the source file.
%% 
%% Any modified versions of this file must be renamed
%% with new filenames distinct from sample-manuscript.tex.
%% 
%% For distribution of the original source see the terms
%% for copying and modification in the file samples.dtx.
%% 
%% This generated file may be distributed as long as the
%% original source files, as listed above, are part of the
%% same distribution. (The sources need not necessarily be
%% in the same archive or directory.)
%%
%%
%% Commands for TeXCount
%TC:macro \cite [option:text,text]
%TC:macro \citep [option:text,text]
%TC:macro \citet [option:text,text]
%TC:envir table 0 1
%TC:envir table* 0 1
%TC:envir tabular [ignore] word
%TC:envir displaymath 0 word
%TC:envir math 0 word
%TC:envir comment 0 0
%%
%% The first command in your LaTeX source must be the \documentclass
%% command.
%%
%% For submission and review of your manuscript please change the
%% command to \documentclass[manuscript, screen, review]{acmart}.
%%
%% When submitting camera ready or to TAPS, please change the command
%% to \documentclass[sigconf]{acmart} or whichever template is required
%% for your publication.
%%
%%
\documentclass[manuscript,screen,review]{acmart}

% acmart 已自带：hyperref, xcolor, amsmath, amssymb, graphicx, booktabs
% 不要重复加载，否则会 option clash
\usepackage{multirow}        % 多行单元格
\usepackage{threeparttable}  % 表内脚注（必须）
\usepackage{tabularx}        % 自适应列宽
\usepackage{array}           % >{...} 列修饰
\usepackage{makecell}        % 单元格内换行 \makecell
% 不要 \usepackage{tcolorbox} / mdframed / soul / ulem，会与 acmart 冲突或留下高亮残迹


\usepackage[UTF8]{ctex}
% \usepackage{amsmath}

% 在 preamble（\begin{document} 之前）：
\theoremstyle{acmdefinition}
\newtheorem{definition}{Definition}
\newtheorem{property}{Property}

\theoremstyle{acmplain}
\newtheorem{theorem}{Theorem}
\newtheorem{proposition}{Proposition}

\usepackage{CJKutf8}
\newcommand{\cnlsh}[1]{\textcolor{black}{\begin{CJK}{UTF8}{gbsn}#1\end{CJK}}}


% ─────────────────────────────────────────────────────────
%  宏包
% ─────────────────────────────────────────────────────────
\usepackage{geometry}
\geometry{a4paper, top=2cm, bottom=2cm, left=2cm, right=2cm}
\usepackage{booktabs}        % \toprule / \midrule / \bottomrule
\usepackage{longtable}       % 可跨页长表
\usepackage{pdflscape}       % landscape 环境
\usepackage{amssymb}         % $\rightleftarrows$（对应 ⇄）

% ─────────────────────────────────────────────────────────
%  自定义列类型：左对齐定宽 p 列
% ─────────────────────────────────────────────────────────
\newcolumntype{L}[1]{>{\raggedright\arraybackslash}p{#1}}

%%
%% \BibTeX command to typeset BibTeX logo in the docs
\AtBeginDocument{%
  \providecommand\BibTeX{{%
    Bib\TeX}}}

%% Rights management information.  This information is sent to you
%% when you complete the rights form.  These commands have SAMPLE
%% values in them; it is your responsibility as an author to replace
%% the commands and values with those provided to you when you
%% complete the rights form.
\setcopyright{acmlicensed}
\copyrightyear{2018}
\acmYear{2018}
\acmDOI{XXXXXXX.XXXXXXX}
%% These commands are for a PROCEEDINGS abstract or paper.
\acmConference[Conference acronym 'XX]{Make sure to enter the correct
  conference title from your rights confirmation email}{June 03--05,
  2018}{Woodstock, NY}
%%
%%  Uncomment \acmBooktitle if the title of the proceedings is different
%%  from ``Proceedings of ...''!
%%
%%\acmBooktitle{Woodstock '18: ACM Symposium on Neural Gaze Detection,
%%  June 03--05, 2018, Woodstock, NY}
\acmISBN{978-1-4503-XXXX-X/2018/06}


%%
%% Submission ID.
%% Use this when submitting an article to a sponsored event. You'll
%% receive a unique submission ID from the organizers
%% of the event, and this ID should be used as the parameter to this command.
%%\acmSubmissionID{123-A56-BU3}

%%
%% For managing citations, it is recommended to use bibliography
%% files in BibTeX format.
%%
%% You can then either use BibTeX with the ACM-Reference-Format style,
%% or BibLaTeX with the acmnumeric or acmauthoryear sytles, that include
%% support for advanced citation of software artefact from the
%% biblatex-software package, also separately available on CTAN.
%%
%% Look at the sample-*-biblatex.tex files for templates showcasing
%% the biblatex styles.
%%

%%
%% The majority of ACM publications use numbered citations and
%% references.  The command \citestyle{authoryear} switches to the
%% "author year" style.
%%
%% If you are preparing content for an event
%% sponsored by ACM SIGGRAPH, you must use the "author year" style of
%% citations and references.
%% Uncommenting
%% the next command will enable that style.
%%\citestyle{acmauthoryear}


%%
%% end of the preamble, start of the body of the document source.
\begin{document}

%%
%% The "title" command has an optional parameter,
%% allowing the author to define a "short title" to be used in page headers.
\title{The Agent Experience Transformation}

%%
%% The "author" command and its associated commands are used to define
%% the authors and their affiliations.
%% Of note is the shared affiliation of the first two authors, and the
%% "authornote" and "authornotemark" commands
%% used to denote shared contribution to the research.
\author{Ben Trovato}
\authornote{Both authors contributed equally to this research.}
\email{trovato@corporation.com}
\orcid{1234-5678-9012}
\author{G.K.M. Tobin}
\authornotemark[1]
\email{webmaster@marysville-ohio.com}
\affiliation{%
  \institution{Institute for Clarity in Documentation}
  \city{Dublin}
  \state{Ohio}
  \country{USA}
}

\author{Lars Th{\o}rv{\"a}ld}
\affiliation{%
  \institution{The Th{\o}rv{\"a}ld Group}
  \city{Hekla}
  \country{Iceland}}
\email{larst@affiliation.org}

\author{Valerie B\'eranger}
\affiliation{%
  \institution{Inria Paris-Rocquencourt}
  \city{Rocquencourt}
  \country{France}
}



%%
%% By default, the full list of authors will be used in the page
%% headers. Often, this list is too long, and will overlap
%% other information printed in the page headers. This command allows
%% the author to define a more concise list
%% of authors' names for this purpose.
\renewcommand{\shortauthors}{Trovato et al.}

%%
%% The abstract is a short summary of the work to be presented in the
%% article.
\begin{abstract}
  A clear and well-documented \LaTeX\ document is presented as an
  article formatted for publication by ACM in a conference proceedings
  or journal publication. Based on the ``acmart'' document class, this
  article presents and explains many of the common variations, as well
  as many of the formatting elements an author may use in the
  preparation of the documentation of their work.
\end{abstract}

%%
%% The code below is generated by the tool at http://dl.acm.org/ccs.cfm.
%% Please copy and paste the code instead of the example below.
%%
\begin{CCSXML}
<ccs2012>
 <concept>
  <concept_id>00000000.0000000.0000000</concept_id>
  <concept_desc>Do Not Use This Code, Generate the Correct Terms for Your Paper</concept_desc>
  <concept_significance>500</concept_significance>
 </concept>
 <concept>
  <concept_id>00000000.00000000.00000000</concept_id>
  <concept_desc>Do Not Use This Code, Generate the Correct Terms for Your Paper</concept_desc>
  <concept_significance>300</concept_significance>
 </concept>
 <concept>
  <concept_id>00000000.00000000.00000000</concept_id>
  <concept_desc>Do Not Use This Code, Generate the Correct Terms for Your Paper</concept_desc>
  <concept_significance>100</concept_significance>
 </concept>
 <concept>
  <concept_id>00000000.00000000.00000000</concept_id>
  <concept_desc>Do Not Use This Code, Generate the Correct Terms for Your Paper</concept_desc>
  <concept_significance>100</concept_significance>
 </concept>
</ccs2012>
\end{CCSXML}

\ccsdesc[500]{Do Not Use This Code~Generate the Correct Terms for Your Paper}
\ccsdesc[300]{Do Not Use This Code~Generate the Correct Terms for Your Paper}
\ccsdesc{Do Not Use This Code~Generate the Correct Terms for Your Paper}
\ccsdesc[100]{Do Not Use This Code~Generate the Correct Terms for Your Paper}

%%
%% Keywords. The author(s) should pick words that accurately describe
%% the work being presented. Separate the keywords with commas.
\keywords{Do, Not, Use, This, Code, Put, the, Correct, Terms, for,
  Your, Paper}

\received{20 February 2007}
\received[revised]{12 March 2009}
\received[accepted]{5 June 2009}

%%
%% This command processes the author and affiliation and title
%% information and builds the first part of the formatted document.
\maketitle


\input{sections/sec1_intro}
% \section{Preliminaries}
% \label{sec:prelim}

% Before entering the discussion of agent experience transformation, this section gives formal definitions of three foundational concepts: agent, agent experience, and experience transformation.

% \subsection{Agent}
% \label{sec:agent}

% This paper defines an agent as a system that performs sequential decision-making in an environment according to task context. At each decision step $t$, the agent perceives the current context $c_t$, produces an action $a_t$, the environment optionally returns an objective outcome $o_t$, and optionally receives an evaluative signal $f_t$.

% An agent's action $a_t$ can simultaneously span natural language reasoning chains, code generation, tool calls, environmental manipulation, and other modalities.

% An agent executing a task follows the loop:

% \[
% c_t \xrightarrow{\pi} a_t \xrightarrow{\mathcal{E}} (o_t) \xrightarrow{} (f_t) \xrightarrow{} c_{t+1}.
% \]

% where $\pi$ denotes the agent's policy, $\mathcal{E}$ denotes the environment, and $(\cdot)$ denotes optional elements. This loop is the production mechanism of experience---each round $(c_t, a_t, o_t, f_t)$ constitutes an atomic experience record, formalized in \S\ref{sec:min_unit}.

% \subsection{Agent Experience}
% \label{sec:agent_exp}

% \subsubsection{Minimum Semantic Unit}
% \label{sec:min_unit}

% The minimum semantic record of agent experience is defined as a modality-agnostic four-tuple:

% \[
% e = (c, a, o, f),
% \]

% where $c$ and $a$ are required elements, and $o$ and $f$ are optional.

% \paragraph{Trajectory.} A complete task execution produces a temporally ordered sequence of experience records, denoted $\tau = (e_1, e_2, \dots, e_T)$, where $T$ is the trajectory length. A collection of multiple trajectories is denoted $\mathcal{D} = \{\tau^{(i)}\}_{i=1}^N$.

% \subsubsection{Experience Carriers}
% \label{sec:carriers}

% The minimum semantic unit $e$ describes the semantic layer of experience---it captures what happened. Experience carriers describe the form in which experience exists within the model architecture, and can be classified by degree of explicitness into three categories: Tokenized, Latent, and Parametric.

% \paragraph{Tokenized Carriers ($\mathcal{C}^T$).} Experience carriers that exist explicitly as discrete token sequences at the input side of the model's forward pass. This covers text tokens, visual patch tokens, action tokens, and serialized structural tokens (e.g., serialized graphs or workflows), and is not limited to natural language text. Tokenized carriers are further divided into two subclasses by degree of formalization:
% \begin{itemize}
%   \item Narrative Tokenized ($\mathcal{C}^T_N$): Weakly formalized carriers organized by natural language or perceptual order, reused through language understanding or multimodal understanding. Typical forms include raw trajectories, reflections, summaries, and screenshots.

%   \item Schematic Tokenized ($\mathcal{C}^T_S$): Strongly formalized carriers organized by syntactic or topological structure, reused through parsing, execution, or graph traversal. Typical forms include workflows, SOPs, and knowledge graphs.
% \end{itemize}

% Raw interaction trajectories themselves are a special case of Narrative Tokenized carriers---they are Narrative carriers at zero abstraction (without any refinement), not raw material that exists independently outside the carrier system. Transformations typically take raw as their starting point.

% \paragraph{Latent Carriers ($\mathcal{C}^L$).} Intermediate representations that exist as continuous vectors or hidden states and participate directly in model attention or hidden-state computation. Latent carriers are not equivalent to discrete tokens (they do not occupy context window space as discrete tokens), nor are they fixed in model weights (they do not alter model parameters); they form a bridging layer between the two. Typical forms include KV cache, prefix cache, learnable soft prompts, and continuous memory tokens.

% \paragraph{Parametric Carriers ($\mathcal{C}^P$).} Experience fixed in neural network weight distributions, fully implicit. Parametric carriers are further divided into two subclasses by functional role:
% \begin{itemize}
%   \item Policy Parameters ($\mathcal{C}^P_\pi$): The actor weights that generate actions $a$.
%   \item Evaluator Parameters ($\mathcal{C}^P_\phi$): The judge weights that evaluate $(c, a, o, f)$.
% \end{itemize}

% \subsection{Experience Transformation}
% \label{sec:transformation}

% Experience transformation refers to the process of moving experience between different carriers. A mapping

% \[
% \mathcal{T}: \mathcal{C}_{\text{src}} \rightarrow \mathcal{C}_{\text{tgt}}
% \]

% constitutes an experience transformation if and only if it jointly satisfies the following two conditions, where $\mathcal{C}_{\text{src}}, \mathcal{C}_{\text{tgt}} \in \{\mathcal{C}^T_N,\, \mathcal{C}^T_S,\, \mathcal{C}^L,\, \mathcal{C}^P_\pi,\, \mathcal{C}^P_\phi\}$:

% \begin{enumerate}
%   \item Experience Grounding. The source content is traceable to one or more experience records $e = (c, a, o, f)$.
%   \item Experience Embodiment. The target carrier encodes the semantic content of the source experience, rather than merely undergoing a change in format.
% \end{enumerate}

\section{Preliminaries}
\label{sec:prelim}

在讨论智能体经验转化之前，本节首先对一些基础概念进行形式化定义：智能体、智能体经验、经验载体和经验转化。

\subsection{Agent}
\label{sec:agent}

本文将 Agent 定义为一个根据任务上下文在环境中进行序列决策的系统:每一步依据当前可获得的信息选择一个动作,环境据此更新,新的信息进入下一步。沿用序贯决策 agent 的标准刻画 [Russell and Norvig 2021],将第 $t$ 步的决策写为

$$a_t \sim \pi_\theta(\cdot \mid c_t),$$

其中 $c_t$ 是该步外部可见的上下文(指令、历史交互、检索内容、环境观测),$\theta$ 是 agent 的内部参数,$\pi_\theta$ 为由参数决定的策略。

这条决策规则把决策依据分置于两处:流经 $c_t$ 的外部信息,与固化在 $\theta$ 中的内部能力，Agent 同时使用两者。上下文窗口中显式可见的交互记录(对话历史、检索文档、工具输出),以及训练所得的权重(世界知识、推理模式、任务技能)。同一份经验既可存在于 $c_t$ 一侧(外部、显式),也可存在于 $\theta$ 一侧(内部、隐式)。

\subsection{Agent Experience}
\label{sec:agent_exp}

Agent 与环境的每一次交互产生一个决策事件,记录为一个模态无关的四元组:

$$e = (c, a, o, f),$$

其中 $c$(Context)为决策前可获得的信息,$a$(Action)为该步输出的动作,$o$(Observation)为环境返回的客观后果,$f$(Feedback)为对 $(c,a,o)$ 或 $(c,a)$ 的评价信号。$c$ 与 $a$ 必需,$o$ 与 $f$ 可选——一步动作未必产生可观测后果,也未必收到评价反馈,缺失不影响记录的完整性。一次完整任务执行产生按时序排列的经验序列,记为轨迹 $\tau = (e_1, \dots, e_T)$;多条轨迹的集合记为 $\mathcal{D} = \{\tau^{(i)}\}_{i=1}^N$。




\subsection{Experience Carriers}
\label{sec:carriers}

最小语义单元 $e$ 描述经验的语义内容;Experience Carrier 描述经验在模型架构中的存在形式,是刻画转化路径源端与目标端的基础坐标。根据经验在模型架构中的存在层次，经验载体可以划分为三类：Tokenized Carrier、Latent Carrier 和 Parametric Carrier。

\paragraph{Tokenized Carrier（$\mathcal{C}^T$）} 以离散 token 序列显式存在于前向传播的输入端,涵盖 text tokens、visual patch tokens、action tokens、序列化的 graph/code/workflow token 等,不限于自然语言。占用 context window,推理时需完整处理。按形式化程度分两子类:
\begin{itemize}
  \item Narrative Tokenized，（$\mathcal{C}^T_N$）：按照自然语言组织的弱形式化载体，通过语言理解或多模态理解实现复用。其典型形式包括原始轨迹、反思、摘要和屏幕截图。
  \item Schematic Tokenized，（$\mathcal{C}^T_S$）：按照句法结构或拓扑结构组织的强形式化载体，通过解析、执行或图遍历实现复用。其典型形式包括工作流、标准操作流程（SOP）和知识图谱。
\end{itemize}

Trajectories 是抽象层级最低的 Narrative Tokenized 载体,转化通常以它为起点。

\paragraph{Latent Carrier（$\mathcal{C}^L$）} 以连续向量或隐藏状态形式存在，并直接参与模型注意力计算或隐藏状态计算的中间表征。潜在载体既不等同于离散词元——它们不会以离散词元的形式占用上下文窗口空间，也不固定在模型权重中——它们不会改变模型参数；因此，潜在载体构成了词元化载体与参数化载体之间的桥接层。其典型形式包括 KV 缓存、前缀缓存、可学习软提示和连续记忆词元。

\paragraph{Parametric Carrier（$\mathcal{C}^P$）} 固化在神经网络权重分布中的经验，完全隐式。推理时不占用 context,直接前向输出,修改需再训练。按功能角色分两子类:
\begin{itemize}
  \item Policy Parameters （$\mathcal{C}^P_\pi$）：生成动作 $a$ 的 actor 权重，涵盖 LLM agent、VLA、GUI agent 权重等。
  \item Evaluator Parameters （$\mathcal{C}^P_\phi$）：评估 $(c,a,o,f)$ 的判断器权重,涵盖 reward models、PRM、verifiers、critics、VLM judges等。
\end{itemize}

\subsection{Experience Transformation}
\label{sec:transformation}

经验转化指经验被重新表示的过程,记为 $\mathcal{T}: \mathcal{C}_{\text{src}} \rightarrow \mathcal{C}_{\text{tgt}}$。重新表示要求目标表示是从源经验重新编码而来的新形式:抽象层级的提升(trajectories → reflections)、形式化程度的改变(trajectories → code)、存在层次的迁移(tokenized → parametric)。源端与目标端是否同属载体不影响判定，比如Narrative → Narrative 的语义抽象同样产生新的经验表示。


\input{sections/sec3_token2token}
\input{sections/sec4_token2latent}
\section{Tokenized-to-Parametric Transformations}

经验从 Tokenized 载体到 Parametric 载体的转化包含两条路径：Evaluator Internalization 把交互经验固化为评估能力，Policy Internalization 把交互经验固化为决策能力。

\subsection{Evaluator Internalization: Tokenized-to-Evaluator Transformation}

Evaluator Internalization 把 Tokenized agent experience 内化为 Evaluator 参数，使评估能力固化在权重中。产物涵盖 reward model（outcome / process）、verifier、critic 与 failure detector 等形态。纳入前提是 Evaluator 参数被经验数据实质更新；纯 prompt 触发、参数不变的 LLM-as-a-judge 不在此列。

\subsubsection{Outcome-supervised Evaluator Internalization}

Outcome-supervised internalization 把监督信号绑定到一条完整 response、整段轨迹或整个 episode，将其映射为单一评价标签。按输出形态，本节分为判别式与生成式 outcome 评估器两支。

\paragraph{Discriminative Outcome Evaluators}

判别式 outcome 评估器输出标量判定，按信号类型分 binary outcome judgment 与 graded outcome scoring 两类。

Binary outcome judgment 将完整 episode 映射为成功/失败或违规/合规等二值判定。Du et al. \cite{Du23} 把 success detection 写成 visual question answering，用 clip-level 成功标签训练 Flamingo 做二值判定。I-FailSense \cite{Gri25} 扩展到跨任务 failure detection，正例来自原指令、负例来自语义不匹配指令。WebRL \cite{Qi24} 在自演化课程 RL 中训练 outcome reward model，依据指令、历史动作与页面状态判定 web trajectory 是否完成任务。Wen et al. \cite{Wen25b} 把判定目标从 task success 移到 policy violation，对四类规则逐条判违规。

Graded outcome scoring 在二值之外给出更细的信号——分级分值、连续概率或相对偏好。RoboReward \cite{Lee26b} 在真实 rollout 上做 counterfactual relabeling 与 negative clipping，构造 1–5 级 rubric 回归离散 progress。SAFE \cite{Gu25b} 以 VLA policy 的内部 latent 为输入，在 rollout prefix 上学习随时间变化的 failure probability，由 conformal threshold 触发 abort。VLP \cite{Liu25w} 改学相对偏好，从 expert / medium / random 三档轨迹构造比较标签，按 Bradley-Terry 学习 cross-modal preference。

\paragraph{Generative Outcome Evaluators}

生成式 outcome 评估器除 verdict 外还产出自然语言文本，按文本承担的角色分 verdict justification、failure diagnosis 与 corrective critique 三类。

Verdict justification 在判定之外生成支撑该判定的推理。GenRM \cite{Zha24o} 把 reward modeling 写成 next-token prediction，有直出 Yes / No 的 GenRM-Direct 与先生成 CoT rationale 再判的 GenRM-CoT。S2J \cite{Sun25l} 让 judge 判断前先自行解题，把 solving 与 judging reward 联合优化。J4R \cite{Xu25i} 让 judge 在 response order swap 等等价输入上保持一致判断，压低位置偏置。AgentV-RL \cite{Zha26ac} 把 verifier 做成多轮 tool-augmented agent，以 verdict 与 ground truth 的一致性为奖励。CriticGPT-RM \cite{Liu24g} 从不同阶段操作视频自动生成 pairwise analysis 与 preference，下游经 Bradley-Terry 训练独立 reward model。

Failure diagnosis 在判定之外生成对失败的解释或定位。AHA \cite{Dua24} 通过 FailGen 在成功 demonstration 上注入多类失败模式，训练模型同时输出 yes / no 与自由文本 failure rationale。ARMOR \cite{Qi26} 联合训练 detection head 与 reasoning decoder，多轮生成 detection 与 reasoning 后按自信度选取结果。ExeVRM \cite{Son26d} 以执行视频训练，输出 video-level success judgment 与失败定位。

Corrective critique 生成供下游 generator 或 actor 取用的修正性意见。CTRL \cite{Xie25c} 以单测通过为奖励用 GRPO 优化 critique，使 generator 更可能产出通过测试的代码。Co-Evolving Critics \cite{Li26l} 让 critic 给出 diagnosis、actor 据此重写轨迹，critic 奖励含 saturation-aware gain 以防随 policy 进化而失效。RL4F \cite{Aky23} 以 refined output 相对 ground truth 的任务指标为奖励训练 critique generator，是 critique-driven repair 的早期形态。

\subsubsection{Process-supervised Evaluator Internalization}

Process-supervised internalization 把监督信号分配到中间步骤或动作前缀，将局部决策质量写进 Evaluator 参数。按输出形态，本节分为判别式与生成式 process 评估器两支。动因来自一项经验观察：长程 agent 的失败通常是若干局部错误的累积，单一 outcome label 无法把这些错误从一条整体失败的轨迹中分离出来。

\paragraph{Discriminative Process Evaluators}

判别式 process 评估器对每个中间步骤或动作前缀输出标量分。按信号语义分 step correctness scoring、step advantage and progress estimation 与 step-action compatibility 三类。

\textbf{Step correctness scoring} 把每步映射为正确/错误或更细的错误类型标签。Let's Verify Step by Step \cite{Lig23} 在 PRM800K 上用人工 step-level 标签训练 PRM，对每一步预测 positive/negative/neutral。OmegaPRM \cite{Luo24} 用 MCTS 与 binary search 自动定位 first error，以 Monte Carlo 估计替代人工标注，产出 soft process label。PathFinder-PRM \cite{Pal25} 的核心主张是 reward 估计应以错误类型为条件——模型先显式判定该步犯了 math error 还是 consistency error，再基于错误类型预测 step reward，而非从步骤文本直接跳到分数。GUI-Shepherd \cite{Che25h} 把 step correctness 标签落到 GUI 域，人工标注 state-action 对是否正确、GPT-4o 补充 CoT rationale 作为训练增强，推理时输出标量正确性分数。GAIA \cite{Wan26p} 把训练组织为 data flywheel——首轮以正确/错误 action 与 ground truth 比对训练 critic，次轮把 critic 难例回流训练第二轮，持续吸收真实 GUI agent 的错误分布。

\textbf{Step advantage and progress estimation} 在二值之外给出更细的信号：估计一步对未来成功概率的边际贡献或任务进度的连续值。Rewarding Progress \cite{Set24} 从每个 prefix 展开 Monte Carlo rollout 估计未来成功概率的变化量，让 Evaluator 直接学习 step advantage 而非 step correctness。PQM \cite{Li24m} 把问题形式化为 Q-value ranking——用 Plackett-Luce 排序损失替代独立 BCE，强制正确前缀的估计成功概率高于错误前缀。BiRM \cite{Che25s} 在同一 backbone 上联合训练 backward-looking reward head 与 forward-looking value head，同时输出步骤局部正确性与 partial trajectory 仍通向正确答案的概率。AgentPRM \cite{Xi25b} 将这一思路从 reasoning 扩展到多步 agent 决策——用 TD 与 GAE 从 rollout 自动估计 state-action 对上的 Q-value 与 advantage，使 agent trajectory 的中间步骤也能获得 process reward。ProgRM \cite{Zha25z} 在 GUI 轨迹中用 LCS recipe matching 自动定位 key step 并分配归一化 progress label，基于动作历史与当前观测预测任务进度。Agentic-Q \cite{Wan26s} 让终局 success 信号沿 trajectory 反传，训练 step-wise success probability estimator。SWE-Shepherd \cite{Dih26} 在 code 域面临一个特殊困难——代码修改轨迹缺少自然的过程监督信号——于是用相关文件访问、目标文件修改、测试结果变化等启发式指标构造弱 process reward，折扣累积为 step-wise target 训练 repo-level PRM。

\textbf{Step-action compatibility} 在二值正确性之外引入另一种信号：比较同一决策点多个候选动作的相对优劣。V-Droid \cite{Dai25} 的核心机制是 Pairwise Process Preference——不孤立地给单个动作打分，而是强制同一步内正确动作的分数高于若干错误动作，从相对序而非绝对值中学习。IntentScore \cite{Che26s} 把 compatibility 的条件从 state 扩展到 planning intent——action encoder 显式接收 planning intent 作为输入，学习 state-action-intent 三元组的相容性分数，使评估器能区分"动作本身正确"与"动作在当前意图下正确"。

\paragraph{Generative Process Evaluators}

生成式 process 评估器在标量判定之外产出自然语言文本——critique、错误诊断或修正建议。按文本在下游的功能分 step-level critique、verdict with structured reasoning 与 pre-operative corrective feedback 三类。

\textbf{Step-level critique} 对中间步骤生成针对性诊断，指出错误位置、类型与原因，不强制附带数值评分。DeepCritic \cite{Yan25u} 先用 teacher model 生成的 deliberate critique 做 critique teaching SFT，再用 PRM800K 与 MC 自动标注的 step correctness 数据做 critique incentivization GRPO，使模型学会多视角、可纠错的长文本步骤批注。AutoMathCritique \cite{Xi24} 让 actor 在 GSM8K 与 MATH 上生成 flawed reasoning paths，由 GPT-4o 标注首个错误步骤，经 Monte Carlo soft filtering 保留真正有助修正的 critique 做 SFT，产物是可直接驱动 actor refinement 的 step-level 反馈。

\textbf{Verdict with structured reasoning} 在给出判定（正确/错误、有希望/无希望）的同时生成支撑推理，使判定可审计、可质疑。与纯 critique 不同，这类评估器需同时产出 verdict 与 reasoning，且两者之间存在逻辑约束——reasoning 必须支撑 verdict。StepWiser \cite{Xio25c} 让模型在对 reasoning chunk 给出 Positive/Negative verdict 前先写出 meta-reasoning，用 MC continuation 估计的 chunk 前后 Q-value 变化作为监督。OS-Oracle \cite{Wu25m} 通过四类错误注入合成负例，用 CP-GRPO 强化 verdict 与 rationale 之间的一致性。WebArbiter \cite{Zha26ab} 在 reasoning 结构上引入一个关键变化：模型在比较候选动作前必须先归纳适用于当前任务与页面的原则，再从原则推导 preference verdict，而非直接从状态跳到判断。Web-Shepherd \cite{Cha25b} 引入另一种结构——checklist-grounded evaluation——按预定义 checklist 逐项生成 Yes/No/In-Progress 判断，使评估过程被分解为可逐项核验的子问题。SOLE-R1 \cite{Sch26} 把 structured reasoning 扩展到具身时序域，从机器人视频中学习生成 per-timestep progress 分数及解释其判断的 spatiotemporal CoT。PPRM \cite{Xu25h} 在不可验证的多轮 search 任务上按 correctness、relevance、consistency 等原则生成 step score，再通过 ReNorm 与 outcome reward 校准。

\textbf{Pre-operative corrective feedback} 在动作执行前生成诊断与修正建议，使错误在发生前被拦截。GUI-Critic-R1 \cite{Wan25q} 在候选动作执行前依次生成 observation analysis、possible result、critique、correctness score 与 corrective suggestion，用 Suggestion-aware GRPO 强化，使 critic 的建议能实际降低后续执行失败率。

% 所需宏包：booktabs, graphicx
% \usepackage{booktabs}
% \usepackage{graphicx}

\begin{table}[htbp]
\centering
\caption{Overview of Evaluator Internalization methods.}
\label{tab:evaluator-internalization}
\resizebox{\textwidth}{!}{%
\begin{tabular}{@{}lllllll@{}}
\toprule
\textbf{Work} & \textbf{Supervision Granularity} & \textbf{Output Form} & \textbf{Evaluator Output} & \textbf{Label Source} & \textbf{Training Method} & \textbf{Domain} \\
\midrule
% outcome / disc
Du23                            & outcome & disc & binary label              & environment-outcome                    & cls-CE          & embodied        \\
I-FailSense~\cite{Gri25}        & outcome & disc & binary label              & ground-truth match                     & cls-CE          & embodied        \\
WebRL~\cite{Qi24}               & outcome & disc & binary label              & environment-outcome                    & cls-CE          & web             \\
Wen25b                          & outcome & disc & binary label              & rule-based                             & cls-CE          & web             \\
RoboReward~\cite{Lee26b}        & outcome & disc & graded score              & environment-outcome                    & regression      & embodied        \\
SAFE~\cite{Gu25b}               & outcome & disc & graded score              & environment-outcome                    & regression      & embodied        \\
VLP~\cite{Liu25w}               & outcome & disc & preference                & environment-outcome                    & ranking         & embodied        \\
\addlinespace[3pt]
% outcome / gen
GenRM~\cite{Zha24o}             & outcome & gen  & structured verdict        & ground-truth match                     & gen-CE          & text/reasoning  \\
S2J~\cite{Sun25l}               & outcome & gen  & structured verdict        & ground-truth match                     & RL-PG           & text/reasoning  \\
J4R~\cite{Xu25i}                & outcome & gen  & structured verdict        & ground-truth match                     & RL-PG           & text/reasoning  \\
AgentV-RL~\cite{Zha26ac}        & outcome & gen  & structured verdict        & ground-truth match                     & gen-CE + RL-PG  & text/reasoning  \\
CriticGPT-RM~\cite{Liu24g}      & outcome & gen  & preference                & environment-outcome$^\dagger$          & gen-CE          & embodied        \\
AHA~\cite{Dua24}                & outcome & gen  & failure diagnosis         & rule-based                             & gen-CE          & embodied        \\
ARMOR~\cite{Qi26}               & outcome & gen  & failure diagnosis         & environment-outcome + human            & cls-CE + gen-CE & embodied        \\
ExeVRM~\cite{Son26d}            & outcome & gen  & failure diagnosis         & environment-outcome                    & gen-CE          & GUI             \\
CTRL~\cite{Xie25c}              & outcome & gen  & corrective suggestion     & environment-outcome                    & RL-PG           & code            \\
ECHO~\cite{Li26l}               & outcome & gen  & corrective suggestion     & environment-outcome                    & RL-PG           & text/reasoning  \\
RL4F~\cite{Aky23}               & outcome & gen  & corrective suggestion     & ground-truth match                     & RL-PG           & text/reasoning  \\
\midrule
% process / disc
Let's Verify~\cite{Lig23}       & process & disc & binary label              & human                                  & cls-CE          & text/reasoning  \\
OmegaPRM~\cite{Luo24}           & process & disc & binary label (soft)       & MC-rollout                             & cls-CE          & text/reasoning  \\
PathFinder-PRM~\cite{Pal25}     & process & disc & binary label (error-typed)& human + MC-rollout                     & cls-CE          & text/reasoning  \\
GUI-Shepherd~\cite{Che25h}      & process & disc & graded score              & human + teacher-model                  & cls-CE          & GUI             \\
GAIA~\cite{Wan26p}              & process & disc & binary label              & ground-truth match                     & cls-CE          & GUI             \\
Rewarding Progress~\cite{Set24} & process & disc & progress/advantage        & MC-rollout                             & regression      & text/reasoning  \\
PQM~\cite{Li24m}                & process & disc & progress/advantage        & MC-rollout                             & ranking         & text/reasoning  \\
BiRM~\cite{Che25s}              & process & disc & progress/advantage        & MC-rollout                             & cls-CE + regression & text/reasoning \\
AgentPRM~\cite{Xi25b}           & process & disc & progress/advantage        & MC-rollout (TD/GAE)                    & regression      & web             \\
ProgRM~\cite{Zha25z}            & process & disc & progress/advantage        & ground-truth match (LCS)               & regression      & GUI             \\
Agentic-Q~\cite{Wan26s}         & process & disc & progress/advantage        & environment-outcome                    & regression      & GUI             \\
SWE-Shepherd~\cite{Dih26}       & process & disc & progress/advantage        & rule-based                             & regression      & code            \\
V-Droid~\cite{Dai25}            & process & disc & preference                & human                                  & ranking         & mobile          \\
IntentScore~\cite{Che26s}       & process & disc & graded score              & ground-truth match                     & ranking         & GUI             \\
\addlinespace[3pt]
% process / gen
DeepCritic~\cite{Yan25u}        & process & gen  & step critique             & teacher-model + MC-rollout             & gen-CE + RL-PG  & text/reasoning  \\
AutoMathCritique~\cite{Xi24}    & process & gen  & step critique             & teacher-model + MC-rollout             & gen-CE          & text/reasoning  \\
StepWiser~\cite{Xio25c}         & process & gen  & structured verdict        & MC-rollout                             & RL-PG           & text/reasoning  \\
OS-Oracle~\cite{Wu25m}          & process & gen  & structured verdict        & rule-based (error injection)           & RL-PG           & GUI             \\
WebArbiter~\cite{Zha26ab}       & process & gen  & structured verdict        & ground-truth match + teacher-model     & gen-CE + RL-PG  & web             \\
Web-Shepherd~\cite{Cha25b}      & process & gen  & structured verdict        & human                                  & gen-CE          & web             \\
SOLE-R1~\cite{Sch26}            & process & gen  & structured verdict        & environment-outcome                    & gen-CE + RL-PG  & embodied        \\
PPRM~\cite{Xu25h}               & process & gen  & structured verdict        & environment-outcome                    & gen-CE          & text/reasoning  \\
GUI-Critic-R1~\cite{Wan25q}     & process & gen  & corrective suggestion     & ground-truth match                     & gen-CE + RL-PG  & GUI             \\
\bottomrule
\end{tabular}%
}
\end{table}

\subsubsection{Discussion}

Outcome label 易于从环境信号、执行结果或人工裁判得到，对具有明确成功标准的任务可直接对齐最终目标，在具身的 episode-level success detection 上几乎是默认选择，下游取用方式直接落在 scalar reward。Process 判别式围绕 credit assignment 快速扩张，演化轴是从 correctness 到 advantage、progress、Q-value、bidirectional reward 的语义细分，再向 GUI、computer-use、code、embodied 等带显式中间状态的领域落地；代价是 step-level label 的获取成本（要么人工标注，要么 MC 自动估计）与噪声。生成式两支让 Evaluator 输出可读 critique 或 diagnosis，下游可塌缩为标量进入 Evaluator-Driven Optimization 的 RL、可取 binary 做 Best-of-N 筛选、也可保留全文驱动 step-wise refinement。生成式评估模型的代价是 critique 的正确性缺乏自动验证，比如一条听起来合理的 critique 可能与环境真实反馈不一致。

跨监督粒度看，process 监督换更细的 credit assignment，付出标注成本与噪声；跨产物形态看，生成式以更高推理成本换更丰富的下游 utilization。两条轴上的选择由下游任务的 verifiability 与监督预算共同决定，不存在普适最优。

\subsection{Policy Internalization: Tokenized-to-Policy Transformation}

Policy Internalization 把 Agent Trajectory，比如自然语言轨迹、工具调用记录、代码编辑历史、GUI 操作序列、reasoning chains， 内化为 Policy 参数，使序贯决策能力固化在权重中。内化后，模型在后续任务中无需检索、拼接或重放原始经验即可决策。按写入 Policy 参数的方式分为三类：imitation-based、environment-reward-based 与 preference-based。

本节的纳入边界由进入 policy update 的监督信号来源决定。信号为非参数化的环境时，比如环境可验证结果、unit test 通过、ground-truth match、rule-based verifier等，归本节；信号来自参数化 evaluator（trained RM / PRM / LLM-as-a-judge）时，不归本节。

\subsubsection{Imitation-based Policy Internalization}

Imitation-based internalization 把高质量 agent trajectories 当作监督示范，通过 behavior cloning、SFT 或 rejection-sampled fine-tuning 把示范动作序列写入 Policy 参数。

主线是 success-filtered behavior cloning。\cite{Pat24} 让 web agent 在 WebArena 中自生成 action-observation trajectories，以环境报错与 self-critique 滤出可信动作，在 QLoRA 下做 autoregressive SFT；SWE-Gym \cite{Pan24} 把过滤信号换成可执行 unit test，以测试通过判定 patch trajectory 成败后做 rejection-sampling fine-tuning。命名有时偏离 objective：AppVLM \cite{Pap25b} 称 RFT / ReST，但最终更新是 success-only MLE，失败轨 return = 0 且不进训练集；具身 VLA 的 \cite{Guo25d} 称 online RL，但 RL 阶段只在 action head 上做探索，最终全模型更新是对 expert 与新发现成功轨并集的 supervised learning。当过滤信号来自参数化 evaluator、但 evaluator 只选数据不进 loss 时，工作仍属 imitation：Go-Browse \cite{Gan25b} 的成功标签由 VLM-as-a-judge 给出、仅对成功轨 SFT，NNetNav \cite{Mur24b} 经 hindsight relabeling 改写交互后用 ORM 过滤轨迹，AndroidGen \cite{Lai25d} 用 StepCritic 切出成功子段再做 LoRA-SFT。

示范来源可来自更强的 teacher。Explorer \cite{Pah25} 用 GPT-4o 多 agent pipeline 合成 multimodal web trajectories，经 Task Verifier 过滤后让学生模型学 next action；AgentTuning \cite{Zen23d} 把多环境 teacher trajectories 与通用 instruction data 混合，使模型在习得 agent 行为协议的同时保持 general instruction-following。

\subsubsection{Environment-Reward-based Policy Internalization}

Environment-Reward-based internalization 让 agent 与环境交互产生 trajectories，用非参数化、可程序验证的反馈直接进入 policy optimization 的 objective。可验证信号涵盖任务成败、unit test pass / fail、execution output、SQL execution correctness、ground-truth answer match、GUI 任务完成状态、机器人操作成败等环境本身可判定的结果；成功与失败轨迹都进入优化，reward 把前者转为正信号、后者转为负信号。困难出现在长程多轮的真实环境：一条数十步的轨迹只在终局拿到单个稀疏 verifiable reward，信号难以利用。相关工作围绕这一困难展开。

% 基本可行性由两项早期工作确立。GLAM \cite{Car23} 把 FLAN-T5 放入 BabyAI-Text、以 sparse goal reward 经 PPO 在线 grounding；LLaRP \cite{Szo23} 连接冻结的 LLaMA 与视觉编码器、以 success / subgoal / invalid-action penalty 组成的环境 reward 经 DD-PPO 训练 adapter 与 heads。动作可执行、环境能返回可验证结果时，Tokenized trajectory 即可经 RL 直接转化为 Policy 参数。

\textbf{Exploration preservation.} 长程 agent 在训练早期极易坍缩到低探索策略，一旦坍缩便不再产生成功轨迹、梯度随之消失。EPO \cite{Xu25k} 在 PPO 或 GRPO 外层加 trajectory-level entropy regularization 与 entropy smoothing，缓解 sparse-reward multi-turn RL 的 exploration-exploitation cascade failure。沿熵调控扩展的有：Agentic Reinforced Policy Optimization \cite{Don25d} 在高不确定性 tool step 触发局部 branching，AEPO \cite{Don25c} 把 entropy 同时纳入 advantage 与 clipping 并动态分配 global 与 branch rollout budget，EMPG \cite{Wan25ad} 用当前 entropy 缩放每步 advantage、再加鼓励进入更低熵后继状态的 future-clarity bonus。另一条思路是保证成功信号不断流：ARPO \cite{Lu25f} 在 GUI 环境维护成功 trajectory replay buffer，采样组全失败、无梯度时从缓冲注回正样本；SoLS \cite{Pap25c} 在 mobile app control 中对正优势样本强更新、对负优势样本做 PPO-like 保守正则，并配合成功转移回放 STR；AgentGym-RL \cite{Xi25c} 的 ScalingInter-RL 用先短后长的 horizon curriculum，让 agent 先在可完成的短程任务上拿到信号。

\textbf{Fine-grained credit assignment.} 归因粒度从整轨持续下沉到 turn、step、token。WebAgent-R1 \cite{Wei25} 以 M-GRPO 并行采样 browser trajectories、把 binary success reward 转成 group-relative advantage；\cite{Gol25} 在 SWE 环境以 unit test pass / fail 加 trajectory length penalty 做 token-level DAPO。actor-critic 一支用 learned value 估计中间 advantage：ArCHer \cite{Zho24f} 先用环境 reward 训练 utterance-level Q critic、再把 multi-turn advantage 传给 token actor，Turn-PPO \cite{Li25ae} 把多轮交互建模为 turn-level MDP、在 turn 粒度用 critic 做 GAE。critic-free 一支从 group 比较中估计相对优势：GiGPO \cite{Fen25c} 用 anchor-state grouping 比较不同 rollout 在相同中间状态下的动作、同时构造 episode-level 与 step-level relative advantage，LOOP \cite{Che25af} 在 AppWorld API 环境用 leave-one-out PPO 做 token-level 更新。另有把信号挂到中间语义的做法：IGPO \cite{Wan25y} 用每轮交互前后 ground-truth answer log-prob 的增量定义 information-gain intrinsic reward，把 observation 的真实信息贡献转成 turn-level dense 信号；RLVMR \cite{Zha25ao} 给轨迹挂 \texttt{<planning>}、\texttt{<explore>}、\texttt{<reflection>}、\texttt{<monitor>} 标签，用 rule-verifiable process reward 加 outcome reward 让 meta-reasoning 行为可被显式监督。

\textbf{Verifiable reward densification.} 改造 reward 信号本身：从环境挖掘可程序验证的中间结构，把单一终局 flag 写成 dense composite reward，使密集信号在每步可得。ToolRL \cite{Qia25b} 把 tool-use reward 拆成 format correctness 与 tool name / arg name / arg value 三层匹配。text-to-SQL 域最集中：SQL-Trail \cite{Hua26d} 在 reason-execute-observe 循环里用 execution、turn-budget、schema grounding 等六项组成 composite reward，SQL-ASTRA \cite{Li26r} 的 Column-Set Matching Reward 从结果表列值集合比较给 dense partial-credit、再用 Asymmetric Trajectory Reward 聚合整轨。LongNav-R1 \cite{Hu26e} 在长程 navigation 中用 geodesic progress shaped reward 加 terminal SPL、以 critic-free time-aware baseline 做 REINFORCE++ 风格更新。GUI 与 VLA 域把 rule-based verifiable reward 直接送入 GRPO：GUI-R1 \cite{Luo25b}、AgentCPM-GUI \cite{Zha25an}、SimpleVLA-RL \cite{Li25aa} 在 format 合法性之上叠加 action-type、click-point、text-match、bbox 命中等可判定项，UI-S1 \cite{Lu25j} 在 semi-online RL 中用模型自生成历史、偏离 expert 时由 patch module 续接，reward 由 offline ground-truth matching 提供。VLM as Decision-Making Agents \cite{Zha24s} 把 CoT 加 action 当作端到端 text rollout，用 λ 重加权 reasoning token 的 log-prob 贡献后以环境 reward 做 PPO。同一范式被扩到多任务与泛化研究：AgentRL \cite{Zha25ag} 以 deterministic verifiable rewards 加 task advantage normalization 稳住多任务异步 RL，Generalization in Mobile Agents \cite{Gu26b} 在异步 emulator 上验证 instance、template、app 三层泛化。Q-SFT \cite{Hon24} 处在本子类的形式边界——形式接近 SFT，但 sample weight 由 Bellman target / Q-value 给出，机制更接近 offline value-based RL，reward provenance 仍是环境信号。

\subsubsection{Preference-based Policy Internalization}

Preference-based internalization 从 agent experience 中构造轨迹间或步骤间的相对偏好，用 DPO 或其它 trajectory-pair optimization 直接更新 Policy。源信号是相对序：优势项来自成功任务、通过验证的执行、修正后的轨迹或更高效的操作路径，劣势项来自失败尝试、错误工具调用或未通过测试的代码修改。偏好的粒度从整轨下沉到轨迹内部的局部单元。

\textbf{Trajectory-level preference.} 偏好定义在整条轨迹上。DMPO \cite{Shi24c} 把 DPO 从单轮 response preference 扩到 multi-turn trajectory preference，引入 state-action occupancy measure 与 length normalization 处理长轨迹下的 preference objective。ETO \cite{Son24} 先用 expert trajectory 做 behavior cloning、再让 policy 主动探索，把成功示范与探索失败轨配成 failure-success pair 持续 DPO。Agent-RLVR \cite{Da25} 在 SWE 环境用 guidance 帮助模型探索更可能成功的修复轨迹，再依 unit test pass / fail 构造 preference pairs 做 DPO——guidance 在此是数据生成器，不进 reward。

\textbf{Sub-trajectory preference.} 偏好下沉到轨迹内部的局部单元。\cite{Che24k} 在 tool-use inference tree 中把"同一父节点下成功 step 优于错误分支 step"作为 step-level preference pair，在 SFT 后接续 DPO。HPL \cite{Gao25c} 同时引入 trajectory、step、action group 三层偏好，group 由 fixed、entropy 或 semantic segmentation 划出，最终目标是 BC 与三层 DPO 的复合，其中 group-DPO 是主要增益来源。WEPO \cite{Liu25z} 把偏好收缩到 web element 层面，用目标元素与 DOM 距离近的干扰元素做对比训练 element selection policy，标签来自 benchmark GT 与 DOM 结构，不依赖 learned judge。

% 所需宏包：booktabs, graphicx
% \usepackage{booktabs}
% \usepackage{graphicx}

\begin{table}[htbp]
\centering
\caption{Overview of Policy Internalization methods.}
\label{tab:policy-internalization}
\resizebox{\textwidth}{!}{%
\begin{tabular}{@{}lllllll@{}}
\toprule
\textbf{Work} & \textbf{Learning Objective} & \textbf{Credit Granularity} & \textbf{Supervision Signal} & \textbf{Optimization Algorithm} & \textbf{Experience Source} & \textbf{Domain} \\
\midrule
% imitation
\cite{Pat24}   & imitation  & trajectory & model-judge        & SFT/BC                              & self          & web                                   \\
\cite{Lai25d}  & imitation  & step       & model-judge        & SFT/BC (LoRA)                       & teacher       & mobile                                \\
\cite{Pap25b}  & imitation  & trajectory & task-outcome       & SFT/BC                              & human / self  & mobile                                \\
\cite{Pah25}   & imitation  & trajectory & model-judge        & SFT/BC                              & teacher       & web                                   \\
\cite{Gan25b}  & imitation  & trajectory & model-judge        & SFT/BC                              & self          & web                                   \\
\cite{Zen23d}  & imitation  & token      & task-outcome       & SFT/BC                              & teacher       & tool-use                              \\
\cite{Mur24b}  & imitation  & trajectory & model-judge        & SFT/BC                              & self          & web                                   \\
\cite{Pan24}   & imitation  & trajectory & unit-test          & rejection-sampling                  & teacher       & code                                  \\
\cite{Guo25d}  & imitation  & trajectory & task-outcome       & SFT/BC                              & human / self  & embodied                              \\
\addlinespace[3pt]
% preference
\cite{Shi24c}  & preference & trajectory & ground-truth match & DPO                                 & teacher       & web / embodied / text-reasoning       \\
\cite{Son24}   & preference & trajectory & task-outcome       & DPO                                 & teacher / self& web / embodied / text-reasoning       \\
\cite{Da25}    & preference & trajectory & unit-test          & DPO                                 & self          & code                                  \\
\cite{Che24k}  & preference & step       & model-judge        & DPO (SFT warm-start)                & teacher       & tool-use                              \\
\cite{Gao25c}  & preference & step       & task-outcome       & DPO (traj/step/group) + BC          & teacher       & embodied / web / SQL                  \\
\cite{Liu25z}  & preference & step       & ground-truth match & DPO                                 & human         & web                                   \\
\midrule
% env-reward
\cite{Car23}   & env-reward & step       & task-outcome       & PPO                                 & self          & embodied                              \\
\cite{Szo23}   & env-reward & step       & shaped-composite   & DD-PPO                              & self          & embodied                              \\
\cite{Zho24f}  & env-reward & token      & task-outcome       & hierarchical actor-critic (ArCHer)  & self          & text-reasoning / web                  \\
\cite{Wei25}   & env-reward & trajectory & task-outcome       & M-GRPO                              & self          & web                                   \\
\cite{Gol25}   & env-reward & trajectory & unit-test          & DAPO                                & self          & code                                  \\
\cite{Luo25b}  & env-reward & step       & ground-truth match & GRPO                                & self          & GUI                                   \\
\cite{Zha25an} & env-reward & step       & ground-truth match & GRPO                                & self          & mobile                                \\
\cite{Li25aa}  & env-reward & trajectory & task-outcome       & GRPO                                & self          & embodied                              \\
\cite{Zha24s}  & env-reward & token      & task-outcome       & PPO                                 & self          & text-reasoning / embodied             \\
\cite{Xi25c}   & env-reward & trajectory & task-outcome       & GRPO (ScalingInter-RL)              & self          & web / embodied / tool-use             \\
\cite{Don25d}  & env-reward & token      & shaped-composite   & ARPO                                & self          & tool-use                              \\
\cite{Zha25ag} & env-reward & token      & rule-verifiable    & GRPO                                & self          & embodied / web / tool-use             \\
\cite{Wan25y}  & env-reward & turn       & ground-truth match & GRPO                                & self          & web                                   \\
\cite{Lu25j}   & env-reward & step       & shaped-composite   & semi-online RL                      & human         & GUI                                   \\
\cite{Xu25k}   & env-reward & trajectory & task-outcome       & EPO (on PPO/GRPO)                   & self          & embodied / text-reasoning             \\
\cite{Hua26d}  & env-reward & trajectory & shaped-composite   & GRPO                                & self          & SQL                                   \\
\cite{Hu26e}   & env-reward & turn       & task-outcome       & HAPO                                & self          & embodied                              \\
\cite{Don25c}  & env-reward & token      & shaped-composite   & AEPO                                & self          & tool-use                              \\
\cite{Wan25ad} & env-reward & step       & shaped-composite   & EMPG (on GRPO/DAPO)                 & self          & web / embodied / tool-use             \\
\cite{Lu25f}   & env-reward & trajectory & rule-verifiable    & GRPO                                & self          & GUI                                   \\
\cite{Fen25c}  & env-reward & step       & task-outcome       & GiGPO                               & self          & embodied / web / tool-use             \\
\cite{Qia25b}  & env-reward & step       & shaped-composite   & GRPO                                & self          & tool-use                              \\
\cite{Gu26b}   & env-reward & trajectory & task-outcome       & GRPO                                & self          & mobile                                \\
\cite{Che25af} & env-reward & token      & task-outcome       & LOOP                                & self          & tool-use                              \\
\cite{Li25ae}  & env-reward & turn       & task-outcome       & Turn-PPO                            & self          & web / text-reasoning                  \\
\cite{Hon24}   & env-reward & token      & task-outcome       & Q-SFT                               & human$^\dagger$ & text-reasoning / web / embodied     \\
\cite{Li26r}   & env-reward & step       & shaped-composite   & GRPO                                & self          & SQL                                   \\
\cite{Zha25ao} & env-reward & step       & rule-verifiable    & GRPO                                & self          & embodied / text-reasoning             \\
\cite{Pap25c}  & env-reward & step       & task-outcome       & SoLS (+STR)                         & self          & mobile                                \\
\bottomrule
\end{tabular}%
}
\end{table}


\subsubsection{Discussion}

三类的差别落在 policy update 如何利用一条轨迹：只取其中的成功动作、取带正负号的全部 reward、还是取轨迹之间的相对序。这一选择同时设定各自的信号承载力、训练稳定性与能力上限，三者并非独立旋钮。Imitation-based 以成功动作序列做最大似然，cold-start 阶段能最快教会模型动作空间、输出协议与基本交互流程；正因监督只来自成功示范，示范覆盖之外的策略就无从经 MLE 习得，有限成功路径会把模型留在表层格式或局部启发式，失败轨被整体弃用。Environment-Reward-based 让 reward 同时携带成功与失败的符号，于是试错能发现示范之外的新策略、能力上限最高；同一个开放试错的设定也意味着 reward 稀疏时 credit assignment 困难、online rollout 昂贵、reward hacking 难防，且 RL 的高方差梯度一旦坍缩或被 hack，所损伤的能力不像编辑一条文本规则那样可逐条回退。Preference-based 取轨迹间的相对序：相对序是 supervised 目标，无需估计高方差 policy gradient，训练比 RL 稳，也把失败轨作为 negative 纳入而非弃用；代价是相对序只约束成对轨迹的高下、不提供绝对信号，效果因此压在 preference pair 的构造质量上，trajectory-level 偏好仍可能缺乏细粒度归因。

三类常被接成同一条 pipeline 的不同阶段：imitation 立起格式与协议，reward optimization 或 preference learning 接手能力突破——它们之间因此更多是互补而非互斥（此类阶段串接的 integration 分析见 §7）。Environment-Reward 与 Preference 沿同一条 credit-assignment 轴下沉——前者向 turn / step / token，后者向 step / element / group——分野在监督形式：绝对 verifiable reward 进 RL objective，相对 pair 进 supervised loss。三类还共享一个 Parametric 载体的属性：经验写进权重后不可局部寻址，新增或修正一类行为需要重新训练，而非像 Tokenized 载体那样逐条编辑。

\section{Parametric-Source Transformations}

以 Parametric 载体为源端的转化共两条:Evaluator-Driven Optimization 把固化在评估器参数中的判断能力转入 Policy 参数,Parametric Externalization 把权重中已内化的决策或评估能力外化为可被其他系统消费的 Tokenized artifacts。两条 pathway 的 source 都是其他 pathway 的产物,而非来自原始交互日志的直接监督。

\subsection{Evaluator-Driven Optimization: Evaluator-to-Policy Transformation}

Evaluator-Driven Optimization 把参数化评估器已习得的评价能力经训练转移进 Policy 参数:Policy 作为 actor 产生行为,评估器对这些行为评判,所得判断作为训练信号更新 Policy 权重。本路径的 source 是 Evaluator parameter 而非被评判的 $(c,a,o)$——被评判的行为只是评估器表达判断的对象,真正被转移的经验早在 Evaluator Internalization 阶段就已压缩进评估器权重,Evaluator parameter 既是那一路径的产物,也是本路径的起点。它与 Policy Internalization 的分界在于学习信号的来源:后者把原始轨迹、人类示范或环境成功信号直接写入权重,前者以评估器对 Policy 行为的判断作为信号;判定上要求 Policy 权重确实因评估器信号而更新,仅在 decoding、reranking、best-of-N、tree search 或 MCTS 中受评估器调控而不更新权重的工作不属此路径。

\subsubsection{Discriminative Outcome Evaluator-to-Policy Transfer}

判别式 outcome 评估器对完整输出或完整 trajectory 给出可比较的标量分数或 chosen/rejected 偏好,policy 把这些信号写入参数。按 policy 吸收信号的方式分为三类:Reward Maximization 把标量当 reward 做 RL 最大化,Preference Contrast 让偏好直接驱动对比式优化,Reward-Ranked Selection 用分数筛样本做监督式写入。

\textbf{Reward Maximization.} 评估器的标量分数充当 reward,policy 用 PPO、GRPO、REINFORCE 等 RL 算法将其最大化。Constitutional AI \cite{Bai22b} 的 RL-CAI 阶段由 constitution 引导的 AI judge 对完整回答给偏好,训练出 preference model 后经 RLAIF 让 policy 最大化其分数。RLAIF \cite{Lee23b} 给出 AI 反馈的两种用法:canonical RLAIF 先用 AI 偏好训 RM,d-RLAIF 直接拿 judge 分数当在线 reward。\cite{Ren26b} 进一步免去 ground-truth label:让 LLM judge 回答 meta-question,以 YES 概率合成 scalar reward 再经 GRPO/CISPO 最大化。\cite{Ack26} 与 Policy Filtration \cite{She24c} 都针对 policy 长期优化 imperfect 评估器引发的 reward hacking:前者在 Dr.GRPO 风格目标上加 gradient regularization 与 reference reset,抑制 policy 对 proxy reward 尖锐高点的利用;后者只放 RM 最可靠的高低分样本进 PPO buffer,缓解 noisy RM 的过拟合。FAPO \cite{Din25b}、\cite{Che25v} 与 OS-Themis \cite{Li26p} 的评估器虽具备步级诊断,但压缩为轨迹标量后才进 RL,落在 outcome 一侧(与 §6.1.3 接壤):\cite{Din25b} 把首个无效步骤的定位折成整条 rollout 的 penalty,\cite{Che25v} 用首错定位算出轨迹的部分正确性标量以激活 all-negative group 的梯度,\cite{Li26p} 的在线 RL 主信号则是 trajectory-level reward。

\textbf{Preference Contrast.} 评估器的 chosen/rejected 偏好不蒸馏成独立 reward,直接进入 DPO、IPO、SLiC 等对比式目标更新 policy。OAIF \cite{Guo24} 是 on-policy 实例:在线 LLM annotator 对当前 policy 采样出的回答对给偏好,直接进 DPO/IPO/SLiC,使监督分布始终跟随 policy、避免 RM 随训练 off-distribution。Self-Rewarding Language Models \cite{Yua24} 与 Meta-Rewarding Language Models \cite{Wu24d} 让评估器与 policy 共享同一组参数:前者让同一模型既作 generator 又作 judge,自评分构造 chosen/rejected 对后以 iterative DPO 回灌;后者再加一层 meta-judge 改进 judge 自身评分质量,形成 actor/judge/meta-judge 共演化。\cite{Wu25x} 与 \cite{Fis24} 都提升偏好优化的稳健性:前者用 Bayesian router 按 query 在多个 RM 间动态选路,把路由后的偏好对送入 online DPO;后者在离线设置下把 RM 的 reward difference 经 L2 distillation 与 forward KL 正则蒸馏进 policy 的 implicit reward,提升 label bias、distribution shift 下的稳定性。

\textbf{Reward-Ranked Selection.} 评估器分数只用于筛选高分样本,再以 SFT 写入 policy,以更低算力逼近 RL 的对齐效果。RAFT \cite{Don23} 用 RM 对同一 prompt 的多条回答打分,保留 best-of-K 做 SFT。BOND \cite{Ses24} 把整个 best-of-N 分布(而不仅是高分样本)用 forward/backward KL 蒸馏进单次采样 policy。BoNBoN \cite{Gui24} 离线刻画 best-of-n 分布,以 SFT-BoN 与 IPO-BoN 同时纳入最优与最差样本。\cite{Gon24b} 的 critic LLM 对整条 trajectory 评分,取 top p\% 作为伪示范配合通用数据迭代 SFT。

Xu24d(Mixture of Judges)用 calibrated RM、LLM judge 与规则约束 judge 的混合信号分别实例化 CRPG、CODPO 与 CRRAFT,正好横跨上述三类;其多评估器去偏价值见 §6.1.5。

\subsubsection{Generative Outcome Evaluator-to-Policy Transfer}

生成式 outcome 评估器对完整输出或完整 trajectory 产出 critique、failure explanation 或 revised response 等自然语言内容,并把这些内容(而非标量)作为 policy 的主要监督。

Constitutional AI \cite{Bai22b} 的 SL-CAI 阶段是早期实例:反馈模型按 constitution 对完整回答生成 critique 与 revision,policy 直接对 revision 做 SFT。ECHO \cite{Li26l} 是更近期的完整实例:与 actor 同步更新的 linguistic critic 对完整 rollout 生成 score-aware critique,actor 据此产出 refined rollout,policy 以 GRPO 在 critique 条件下优化这些 refined trajectory,critic 在同批数据上经双轨 GRPO 共演化。生成式反馈的价值天然倾向落到步级——评估器一旦愿意生成细粒度解释,用在 process 粒度比 outcome 粒度信息利用率更高,因此 outcome 形态的生成式评估器多向 §6.1.4 迁移。

\subsubsection{Discriminative Process Evaluator-to-Policy Transfer}

判别式 process 评估器对中间步骤、局部动作、reasoning prefix 或 state-action pair 给出可比较的步级信号,policy 据此更新参数。核心动机是长程 agent 的 credit assignment:仅靠末端 reward,policy 难以辨别该强化或避免哪些局部行为。按步级信号的产生方式分为三类:Step-Value Critic 训练 critic 估计步级 value,Step-Quality Scoring 让评估器直接给每步打质量分,Reward Redistribution 把已有的轨迹级 reward 摊到各步。

\textbf{Step-Value Critic.} 训练一个专门的 critic 估计每步的 value 或 advantage,作为 credit assignment 的 baseline。SWEET-RL \cite{Zho25p} 在 multi-turn 协作推理中训练 turn-level critic,借训练期特权信息直接学每轮 action 的 advantage,再经 DPO 把 turn-wise credit assignment 蒸馏进 actor。\cite{Wan26s} 在 GUI 场景训练 Agentic-Q model 估计 state-action 的 future return,把 step-wise value 接入 SWPO 式更新。

\textbf{Step-Quality Scoring.} 评估器直接对每步或每个维度给出质量分,policy 通过 RL reward、筛选或偏好写入。PPRM \cite{Xu25h} 面向 non-verifiable 的 multi-turn search,用 principle-based PRM 对每轮的 correctness、relevance、coherence 评分,经 ReNorm 与末端 ORM 校准后送入 PPO+GAE。Language Feedback Model \cite{Zho24e} 与 \cite{Pan24b} 都把每步的质量判断转成筛选信号:前者判断每个动作是否 productive、对被标为 desirable 的 state-action 做 behavior cloning;后者用 per-step progress filtering 的结果做 Filtered BC。AdaRubric \cite{Din26e} 与 \cite{Luo25f} 都把质量分构造成偏好:前者先为任务生成 rubric、再对每步每维输出 1–5 分,聚合后做 trajectory-level DPO 与 step-level PPO;后者用 judge 对 computer-use 单步 GUI 动作打 0–100 分,由高低分构造 step-level 偏好对训练 UI-TARS-2B。

\textbf{Reward Redistribution.} 不额外构造步级标签,把已有的轨迹级 reward 摊到 token 或 round 上。\cite{Cha24b} 取标准 RM 最后层 attention 的 token 归因,将终局标量 reward 重分配为 token-level dense reward 送入 PPO。CW-GRPO \cite{Wan26ac} 用 rubric-based judge 对每轮 search 给 retrieval utility 与 reasoning correctness 两个 binary 信号,合成 round-level contribution weight 把 trajectory-level advantage 按轮重分配。

\subsubsection{Generative Process Evaluator-to-Policy Transfer}

生成式 process 评估器对中间步骤产出解释性、诊断性或修正性内容,并把这些内容转进 policy。与判别式 process 回答"这一步好不好"不同,生成式 process 回答"这一步为什么不好、应如何改"。按生成内容的形态分为两类:Step Critique-and-Refine 对每步 critique 并产出 refined step,Failure-Driven Revision 定位错误或给出指令再把修正内化。

\textbf{Step Critique-and-Refine.} 评估器对每步 critique,actor 据此产出 refined step,refined 内容作为对齐目标写入 policy。SRR-Judge \cite{Zha26am} 对 search agent 每个 thought-action step 同时输出 explanation、1–5 rating 与 refined thought/action,rating 负责筛选、refined step 作为生成式对齐目标,经 iterative RFT 写入 policy。Natural Language Actor-Critic \cite{Hon25c} 用 generative language critic 为每个 state-action 产出优劣理由与未来走向的自然语言 critique,actor 在 critique 条件下生成 refined action,再反向蒸馏回基础 policy。

\textbf{Failure-Driven Revision.} 评估器定位错误或给出修正指令,修正后的轨迹或动作经内化写入 policy。Agent-R \cite{Yua25c} 让 actor 自身充当 verifier,在搜索树分叉中定位坏轨迹的 transition point,在错误前缀与正确后续之间插入 revision signal,把"反思加修正"实例做迭代 SFT。OpenClaw-RL \cite{Wan26ad} 从用户话语或环境信号抽取 directive textual hints,用 On-Policy Distillation 把 hint-enhanced teacher 与 student 的 log-prob gap 转成 token-level directional update;其 sequence-level binary RL 旁支不在本格讨论。

% 所需宏包：booktabs, graphicx
% \usepackage{booktabs}
% \usepackage{graphicx}

\begin{table}[htbp]
\centering
\caption{Overview of Evaluator-Driven Optimization methods.}
\label{tab:evaluator-driven-optimization}
\resizebox{\textwidth}{!}{%
\begin{tabular}{@{}llllll@{}}
\toprule
\textbf{Work} & \textbf{Evaluator Granularity} & \textbf{Evaluator Form} & \textbf{Transfer Mechanism} & \textbf{Evaluator Coupling} & \textbf{Domain} \\
\midrule
% outcome / discriminative
% \cite{Bai22b}  & outcome & discriminative  generative & reward-maximization / critique-conditioned-refine    & frozen-separate  & dialogue            \\
\cite{Lee23b}  & outcome & discriminative              & reward-maximization                                  & frozen-separate  & dialogue            \\
\cite{Ren26b}  & outcome & discriminative              & reward-maximization                                  & frozen-separate  & text/reasoning      \\
\cite{Ack26}   & outcome & discriminative              & reward-maximization                                  & frozen-separate  & text/reasoning      \\
\cite{She24c}  & outcome & discriminative              & reward-maximization                                  & frozen-separate  & code                \\
\cite{Din25b}  & outcome & discriminative              & reward-maximization                                  & frozen-separate  & text/reasoning      \\
\cite{Che25v}  & outcome & discriminative              & reward-maximization                                  & frozen-separate  & text/reasoning      \\
\cite{Li26p}   & outcome & discriminative              & reward-maximization                                  & multi-evaluator  & GUI                 \\
% \addlinespace[3pt]
\cite{Guo24}   & outcome & discriminative              & preference-contrast                                  & frozen-separate  & dialogue            \\
\cite{Yua24}   & outcome & discriminative              & preference-contrast                                  & shared-params    & text/reasoning      \\
\cite{Wu24d}   & outcome & discriminative              & preference-contrast                                  & shared-params    & text/reasoning      \\
\cite{Wu25x}   & outcome & discriminative              & preference-contrast                                  & multi-evaluator  & text/reasoning      \\
\cite{Fis24}   & outcome & discriminative              & preference-contrast                                  & multi-evaluator  & text/reasoning      \\
% \addlinespace[3pt]
\cite{Don23}   & outcome & discriminative              & reward-ranked-selection                              & frozen-separate  & dialogue            \\
\cite{Ses24}   & outcome & discriminative              & reward-ranked-selection                              & frozen-separate  & text/reasoning      \\
\cite{Gui24}   & outcome & discriminative              & reward-ranked-selection                              & frozen-separate  & text/reasoning      \\
\cite{Gon24b}  & outcome & discriminative              & reward-ranked-selection                              & frozen-separate  & tool-use            \\
% \cite{Xu24d}   & outcome & discriminative              & reward-maximization / preference-contrast / reward-ranked-selection & multi-evaluator & dialogue \\
% \midrule
% process / discriminative
\cite{Zho25p}  & process & discriminative              & preference-contrast                                  & frozen-separate  & code                \\
\cite{Wan26s}  & process & discriminative              & reward-maximization                                  & frozen-separate  & GUI                 \\
\cite{Xu25h}   & process & discriminative              & reward-maximization                                  & frozen-separate  & search              \\
\cite{Zho24e}  & process & discriminative              & reward-ranked-selection                              & frozen-separate  & embodied            \\
\cite{Pan24b}  & process & discriminative              & reward-ranked-selection                              & frozen-separate  & web / mobile        \\
% \cite{Din26e}  & process & discriminative              & preference-contrast / reward-maximization            & frozen-separate  & web                 \\
\cite{Luo25f}  & process & discriminative              & preference-contrast                                  & frozen-separate  & GUI                 \\
\cite{Cha24b}  & process & discriminative              & reward-redistribution                                & frozen-separate  & text/reasoning      \\
\cite{Wan26ac} & process & discriminative              & reward-redistribution                                & frozen-separate  & search              \\
% process / generative
\cite{Zha26am} & process & generative                  & critique-conditioned-refine                          & frozen-separate  & search              \\
% \cite{Hon25c}  & process & generative                  & critique-conditioned-refine                          & shared-params    & text/reasoning / web / dialogue / tool-use \\
\cite{Hon25c}  & process & generative                  & critique-conditioned-refine                          & shared-params    & general \\
\cite{Yua25c}  & process & generative                  & failure-driven-revision                              & shared-params    & web / embodied      \\
% \cite{Wan26ad} & process & generative                  & failure-driven-revision                              & frozen-separate  & dialogue / GUI / code / tool-use \\
\cite{Wan26ad} & process & generative                  & failure-driven-revision                              & frozen-separate  & general \\
\bottomrule
\end{tabular}%
}
\end{table}

\subsubsection{Discussion}

Evaluator-Driven Optimization 的共性是在经验与 policy 之间插入一个学得的评估器,本路径的能力与局限都由这层中介决定。它让对齐得以规模化——"什么样的输出更好"被离线压进评估器权重后,采样、筛选、对比的边际成本极低,这是 RLHF 工具链横向扩展的前提——并把可训练范围延伸到 rule-based verifier 不及的 non-verifiable 区域:开放对话、长程 search、多维 agentic 行为中环境给不出 clean reward,评估器几乎是把难以量化的偏好转成可优化信号的唯一手段。同一层中介也设下上界,policy 的能力不超过评估器;对一个固定且 imperfect 的评估器持续优化,policy 逼近的是评估器认定的好而非真实目标,reward hacking 因此是本路径的内生风险,在 scalar/critique、outcome/step 各形态下都会复现。

本路径内部沿两条轴展开,各自承载一组张力。粒度轴(outcome / process)决定 credit assignment 的精度:outcome 信号简单,long-horizon credit 却粗糙;process 信号提供 dense supervision,step-label 的可靠性与一致性却更难保证。形态轴(discriminative / generative)决定信息密度与可验证性之间的折中:判别式信号密度低,但可比较、易筛易排;生成式信号密度高,却难以核验,且与 refined trajectory 等 tokenized 载体边界模糊,常令方法在归属上贴近 Policy Internalization 或 Parametric Externalization。这两轴上的演化方向一致——从判别走向生成、从 outcome 走向 process,信息密度随之抬升,而"评估器自身是否可靠"也被推到更前。

评估器的可靠性是本路径的约束核心,也是最大的开放问题。现有缓解各有侧重——co-evolution 让评估器随 policy 更新 \cite{Yua24, Wu24d, Li26l}、ensembling 与 routing 用多评估器分散单点偏差 \cite{Wu25x, Xu24d}、regularization 与 filtration 约束优化过程 \cite{Ack26, She24c}——但都只能压制、无法消除评估器与真实目标之间的缺口;可验证性还呈不对称:判别式信号尚可用 held-out 偏好一致性间接核验,生成式步级 critique 的正确性至今依赖人工或更强模型,而错误的 critique 比错误的 scalar 更具误导性。把本路径放回相邻路径中看,Policy Internalization 的监督直接 grounded 却把噪声一并传入,Evaluator-Driven Optimization 以一层 grounding 换取更稳、更可比的信号,二者在环境信号充足与稀疏时各有所长;评估器本身又是 Evaluator Internalization 的产物,P4 → P6 的串接构成完整的 evaluator-based pipeline,其复合形态与跨路径的系统比较分别留给 §7 与 §8。

\subsection{Parametric Externalization: Parametric-to-Tokenized Transformation}

Parametric Externalization 把已内化在 Policy 或 Evaluator 参数中的经验外化为可被其他系统使用的 Tokenized artifacts。外化产物分两类,对应两端参数的角色:Policy 端产出 agent trajectory(决策行为的完整记录),Evaluator 端产出 evaluative annotation(对行为质量的判断,含 correctness/quality label、自然语言 critique、preference annotation)。与其余路径处理一条已存在的 trajectory 不同,本路径的 artifact 由承载参数的模型现场生成——同一组参数在不同 prompt 与环境下可产出无数条轨迹,转化算子因此是从 $p(e \mid \text{prompt}, \text{env})$ 中采样,而非对某个给定 $e$ 的重编码。

按外化产物来自哪一端参数,本路径分为两类:Demonstration Externalization(§6.2.1)外化 Policy 的决策行为,Evaluative Supervision Externalization(§6.2.2)外化 Evaluator 的判断能力。

\subsubsection{Demonstration Externalization}

Demonstration Externalization 把参数化 Policy 或 teacher agent 的隐式决策能力外化为可供 student 模仿的 agent trajectory,下游以 SFT、behavioral cloning 或 imitation 训练新模型。外化轨迹的真实程度取决于 teacher 生成时所处的环境:teacher 在真实环境中交互时,observation 是环境的客观返回,轨迹 grounding 最强;环境若由模型模拟乃至完全省去,observation 转由模型自身生成,轨迹愈接近一次纯参数读出,pseudo-success 与 hallucinated observation 的风险随之上升,对生成阶段验证的依赖也相应加重。据此分为 Real-Environment Synthesis、Simulated-Environment Synthesis 与 Constructive Synthesis 三类。

\textbf{Real-Environment Synthesis.} Teacher 在真实可交互环境中执行任务,observation 来自环境的客观返回,grounding 最强。web 域的代表是 Explorer \cite{Pah25}:GPT-4o 在真实网页中自主发现任务、执行多步交互,经多阶段 refine 与 verifier 产出带页面状态、动作序列与总结的 multimodal trajectory。沿此扩展,InSTA \cite{Tra25} 推到 internet scale,AutoSurfer \cite{Fai26} 先识别站点功能空间再整理为 task tuple,SynthAgent \cite{Wan25ar} 做 element-level exploration 后全局 refine,A3-SYNTH \cite{Lu26} 以 Agent-as-Annotators 产出带结构化 reasoning 的成功轨迹;\cite{Log26b} 用 judge 抽取 constraints 并按满足率做 prefix extraction 与 hindsight relabeling,使部分成功轨迹也可训练。计算机与移动端,Fara-7B \cite{Awa25} 经 FaraGen 产出 screenshot-thought-action 轨迹并由三重 verifier 筛选,OpenMobile \cite{Che26d} 由 expert 接管修正 learner 偏差、保留 corrected trajectory 与 step-level CoT,AgentSynth \cite{Xie25e} 在 OSWorld 等环境先生成并验证简单 subtask、再合成更长程复合任务并保留完整 action trajectory 与视觉证据。工具与 MCP 域,TOUCAN \cite{Xu25o} 在近 500 个真实 MCP 环境中合成任务并执行完整 tool-use rollout,经 task 级与 trajectory 级双重过滤形成 1.5M 规模 corpus。ANCHOR \cite{Wei26b} 沿少量 high-quality seed trajectory 找 branch point、在分叉状态提新任务并续行,形成共享前缀、后段分叉的树状 GUI 数据。

\textbf{Simulated-Environment Synthesis.} 环境本身由模型扮演,observation 由 simulator 生成。这组在工具调用域最集中:ToolAlpaca \cite{Tan23} 用 ChatGPT 分饰 user、assistant 与 tool executor 模拟生成工具轨迹,ToolMind \cite{Yan25ab} 以三类 agent 模拟产出带 reasoning 的多轮轨迹并做细粒度过滤,ToolACE \cite{Liu24o} 先扩充 API pool 再合成多复杂度的工具调用对话,APIGen-MT \cite{Pra25b} 先定 task blueprint 与 ground-truth action、再经 simulated agent-human interplay 合成多轮 function-calling 对话。mobile GUI 域,Learning with Challenges \cite{Kan26b} 先对 student 做能力 profiling、再用 explorer/supervisor/synthesizer 等 VLM 角色在 emulator 中生成难度自适应的 goal-conditioned trajectory。一组工作把"环境本体"作为主产物:\cite{Wan25as} 把 GPT-4o-mini 同时用作 world simulator 与 teacher、既产 trajectory 又产可扩展的 UI state-transition simulator,Mock Worlds \cite{Lu26j} 从 persona 合成 workflow、虚拟 tool ecosystem 与 rubric,ASTRA \cite{Tia26} 一线产 SFT trajectory、一线把工具实现编译成 sandbox-verifiable RL arena,\cite{Xu26e} 在少量 seed 上 self-evolving 生成后扰动 query 与 tool 构造 RL-ready mock 环境。

\textbf{Constructive Synthesis.} 此端几乎没有交互环境,trajectory 围绕预先给定的 answer、evidence 或问题结构反向构造,最接近一次纯参数读出。EviPath \cite{Li25at} 从 gold evidence 反推 abductive reasoning path 再生成含 planning 与 evidence-grounded answering 的 RAG trajectory,RAGShaper \cite{Tao26} 构造带 distractor 的信息树与多跳问题、在受控检索环境诱导 teacher 产出 agentic RAG trajectory,APIGen \cite{Liu24n} 基于真实可执行 API 生成 query 与 function-calling answer 并经三层检查筛选(保留 execution check,但交互浅、整体偏生成)。

\subsubsection{Evaluative Supervision Externalization}

Evaluative Supervision Externalization 把参数化 Evaluator(judge、verifier、critic、reward model)的隐式评估能力外化为 Tokenized evaluative annotation,用于训练 student、蒸馏轻量 evaluator 或筛选数据。产物按形态分三类:step 或 outcome 层的 correctness/quality label、自然语言 critique 与 failure diagnosis、preference annotation。

\cite{He25g} 先用强 computer-use agent 产出 noisy trajectory,再由 o4-mini 对每步做 correctness 与 optimality grading、附 screenshot/action/alternative 分析,结果被物化为 WebSTAR 与 WebSCORE 两套独立数据,前者训练 student agent,后者蒸馏轻量 StepRM。Step-GUI \cite{Yan25x} 先以轨迹级 verifier 建立质量锚点,再由强多模态 model 把验证结果转写为 progress tracking、state summary、effect prediction 等七类结构化 step-level supervision,与原轨迹一道用于 cold-start fine-tuning 与 RLVR。

% 所需宏包：booktabs, graphicx
% \usepackage{booktabs}
% \usepackage{graphicx}

\begin{table}[htbp]
\centering
\caption{Overview of Parametric Externalization methods.}
\label{tab:parametric-externalization}
\resizebox{\textwidth}{!}{%
\begin{tabular}{@{}llllllll@{}}
\toprule
\textbf{Work} & \textbf{Source Parameter} & \textbf{Environment Grounding} & \textbf{Externalized Artifact} & \textbf{Verification Mechanism} & \textbf{Generator--Learner Relation} & \textbf{Downstream Use} & \textbf{Domain} \\
\midrule
% policy / real-environment
\cite{Pah25}                        & policy    & real-environment       & agent trajectory                  & LLM/VLM-judge          & teacher$\to$student (cross-model) & SFT/BC     & web             \\
\cite{Tra25}                        & policy    & real-environment       & agent trajectory                  & LLM/VLM-judge          & teacher$\to$student (cross-model) & SFT/BC     & web             \\
\cite{Tao26}                        & policy    & constructive           & agent trajectory                  & env-success            & teacher$\to$student (cross-model) & SFT/BC     & RAG             \\
\cite{Xu25o}                        & policy    & real-environment       & agent trajectory                  & multi-stage-verifier   & teacher$\to$student (cross-model) & SFT/BC     & tool-use/MCP    \\
\cite{Li25at}                       & policy    & constructive           & agent trajectory                  & none                   & teacher$\to$student (cross-model) & SFT/BC     & RAG             \\
\cite{Wan25ar}                      & policy    & real-environment       & agent trajectory                  & multi-stage-verifier   & teacher$\to$student (cross-model) & SFT/BC     & web             \\
\cite{Awa25}                        & policy    & real-environment       & agent trajectory                  & multi-stage-verifier   & teacher$\to$student (cross-model) & SFT/BC     & computer-use    \\
\cite{Yan25ab}                      & policy    & simulated-environment  & agent trajectory                  & multi-stage-verifier   & teacher$\to$student (cross-model) & SFT/BC     & tool-use/MCP    \\
\cite{Che26d}                       & policy    & real-environment       & agent trajectory                  & env-success            & teacher$\to$student (cross-model) & SFT/BC     & mobile          \\
\cite{Pra25b}                       & policy    & simulated-environment  & agent trajectory                  & multi-stage-verifier   & teacher$\to$student (cross-model) & SFT/BC     & tool-use/MCP    \\
\cite{Tan23}                        & policy    & simulated-environment  & agent trajectory                  & none                   & teacher$\to$student (cross-model) & SFT/BC     & tool-use/MCP    \\
\cite{Liu24n}                       & policy    & constructive           & agent trajectory                  & multi-stage-verifier   & teacher$\to$student (cross-model) & SFT/BC     & tool-use/MCP    \\
\cite{Liu24o}                       & policy    & simulated-environment  & agent trajectory                  & multi-stage-verifier   & teacher$\to$student (cross-model) & SFT/BC     & tool-use/MCP    \\
\cite{Fai26}                        & policy    & real-environment       & agent trajectory                  & multi-stage-verifier   & teacher$\to$student (cross-model) & SFT/BC     & web             \\
\cite{He25g}            & policy    & real-environment       & agent trajectory                  & multi-stage-verifier   & teacher$\to$student (cross-model) & SFT/BC     & computer-use    \\
\cite{Wei26b}                       & policy    & real-environment       & agent trajectory                  & multi-stage-verifier   & teacher$\to$student (cross-model) & SFT/BC     & GUI             \\
\cite{He25g}    & evaluator & annotation-over-trajectory & correctness/quality label    & rubric                 & teacher$\to$student (cross-model) & RM-training & computer-use   \\
\bottomrule
\end{tabular}%
}
\end{table}


\subsubsection{Discussion}

Parametric Externalization 的价值在两个层面。其一是数据生产:它以 teacher 的自动生成替代人工逐条标注,既降低单位成本,又能沿 curriculum、difficulty、domain coverage、failure pattern 等维度定向合成所需经验,数据规模与覆盖面不再受人工吞吐约束。其二是能力转移:昂贵的大 teacher 的隐式能力可被外化、再蒸馏进一组轻量、领域专用的 student,把推理期持续的高成本换成一次性的合成成本。

代价大多从一个根本约束派生:student 在任务能力维度继承的是 teacher 的上界,本路径兑换的是成本、推理效率与领域专用化,而非能力前沿本身。teacher hallucination、伪成功轨迹与不可执行步骤会被 student 当作正确监督吸收,这一风险沿环境合成度轴递增——真实环境一端 verifier 校验任务成功即可,模拟与构造一端还须校验 observation 自身的可信度;distribution narrowing 则放大单一 teacher 的风格与偏好,multi-teacher mixing 可缓解但不能根除。verifier 因此既是本路径的支柱也是最薄弱处:false positive 让伪成功漏入 student、false negative 错杀高质量样本,如何把 evaluation-grade verifier 与生成规模同步扩展,是本路径的开放工程问题。

在经验来源维度,Parametric Externalization 是 \{teacher\}-sourced experience 进入训练语料的主要途径——其余路径处理的多是 agent 自身或人类产生的经验,而 teacher 模型的隐式能力须经本路径外化,才能成为可供其他系统使用的显式经验。

\input{sections/sec7_composite_pipelines}
\section{Pathway-Level Analysis}

§3–§7 按转化路径逐一梳理了社区将 agent experience 从一种载体重新表示为另一种载体的机制。本章横向比较这些路径与它们产出的载体：每条 pathway 在生产侧付出什么，每类 carrier 在利用侧表现如何，任务约束如何在二者之间做出选择。

Section~\ref{sec:agent} 将 agent 决策写为 $a_t \sim \pi_\theta(\cdot \mid c_t)$，经验影响决策只有两条途径——流经 $c_t$ 的外部信息，或固化在 $\theta$ 的内部能力。Section~\ref{sec:carriers} 据此将载体分为 Tokenized， Latent 和 Parametric：输入侧的离散 token、输入侧的连续向量、权重内的隐式表示。

Section~\ref{sec:prod_profiles} 以五个维度刻画每条 pathway 的生产画像，Section~\ref{sec:util_profiles} 以四个维度刻画每类 carrier 的利用画像，Section~\ref{sec:selection_constraints} 将任务约束读为这两组维度的对偶——分析约束如何剪除不可行载体、约束冲突时如何迫使载体组合、约束随部署漂移时最优配置如何随之改变。

\subsection{Pathway Production Profiles}
\label{sec:prod_profiles}

本节以五个维度刻画每条 pathway 的生产画像。Production Mechanism（D1）：转化靠什么计算机制执行，如 LLM inference、gradient-based training、gradient-free persistence、execution + verification loop。Construction Cost（D2）：转化消耗多少资源。Information Transformation（D3）：经验信息在转化中被保留、损失、注入了什么，以及能否被追溯（traceability）；其中 loss 区分 lossy-by-design（抽象/压缩必然丢弃细节，损失类别可控）与 lossy-by-limitation（欠拟合、reward hacking、采样偏差，损失类别不可控，属工程风险而非设计意图）。Incrementality（D4）：新经验持续到来时，转化能否增量执行。Output Verifiability（D5）：转化产物能否在不依赖下游任务表现的情况下被独立验证。





\begin{landscape}

{\small                          % 字号调小，适配宽表
\setlength{\tabcolsep}{5pt}      % 列间距微调

\begin{longtable}{%
  L{2.0cm}   % Pathway
  L{2.8cm}   % Mechanism
  L{2.5cm}   % Construction Cost
  L{6.5cm}   % Information Transformation
  L{3.5cm}   % Incrementality
  L{2.8cm}   % Output Verifiability
  L{2.2cm}   % 代表工作 / 详节
}

%── 标题（如处于第 8 章，LaTeX 会自动编号为 Table 8.1）──────
\caption{Pathway production profiles across five dimensions.}%
\label{tab:pathway-profiles}\\

%── 首页页眉 ──────────────────────────────────────────────
\toprule
\textbf{Pathway} &
\textbf{Mechanism} &
\textbf{Construction Cost} &
\textbf{Information Transformation} &
\textbf{Incrementality} &
\textbf{Output Verifiability} &
\textbf{代表工作 / 详节} \\
\midrule
\endfirsthead

%── 续页页眉 ──────────────────────────────────────────────
\multicolumn{7}{c}{\tablename~\thetable\quad（续表）}\\[3pt]
\toprule
\textbf{Pathway} &
\textbf{Mechanism} &
\textbf{Construction Cost} &
\textbf{Information Transformation} &
\textbf{Incrementality} &
\textbf{Output Verifiability} &
\textbf{代表工作 / 详节} \\
\midrule
\endhead

%── 续页页脚 ──────────────────────────────────────────────
\midrule
\multicolumn{7}{r}{\small 接下页} \\
\endfoot

%── 末页页脚 ──────────────────────────────────────────────
\bottomrule
\endlastfoot

%── P1 ───────────────────────────────────────────────────
P1 Narrative Abstraction &
LLM inference &
低·可全自动 &
\textbf{保留}: 模式与洞察 /
\textbf{损失}: 步骤与时序(by-design) /
\textbf{注入}: 跨轨迹综合 insight(源自模型先验) &
增量(append); dynamic store 有 merge/prune 级联 &
直接阅读判断(主观) &
Reflexion [Shi23b]; \S3.1 \\
\midrule
%── P2 ───────────────────────────────────────────────────
P2 Schematic Formalization &
LLM inference + 可选 execution 验证 &
中·需可执行环境 &
\textbf{保留}: 过程结构 /
\textbf{损失}: NL nuance/隐含假设(by-design) /
\textbf{注入}: 显式依赖结构 &
增量(新 artifact 入库) &
execution test(客观, coverage-limited) &
Voyager [Wan23c]; \S3.2 \\
\midrule
%── P3 ───────────────────────────────────────────────────
P3 Latent-Space Transformation &
cache: gradient-free / trained: gradient &
cache 低·仅 session; trained 中 &
\textbf{保留}: 计算状态 /
\textbf{损失}: 符号可读性(by-design) /
\textbf{注入}: token 未显式编码的统计规律 &
cache 即时重建; trained 需 batch 重训 &
无直接验证·仅 ablation 间接 &
LAG [Che25b](cache)/
ReasonCACHE [Gup26](trained); \S4 \\
\midrule
%── P4 ───────────────────────────────────────────────────
P4 Evaluator Internalization &
gradient training &
高·需标注或明确成败信号 &
\textbf{保留}: 评估偏好 /
\textbf{损失}: 校准来源与单条贡献
(mix: by-design + by-limitation) &
batch 重训; 新标注需重训以保校准 &
无直接验证·hold-out calibration check &
Let's Verify Step by Step [Lig23]; \S5.1 \\
\midrule
%── P5 ───────────────────────────────────────────────────
P5 Policy Internalization &
gradient training (SFT/RL) &
高·需大规模高质量 trajectory &
\textbf{保留}: 决策模式 /
\textbf{损失}: 单条贡献不可区分(by-design)
+ 欠拟合/reward hacking 风险(by-limitation) &
batch 重训(可热启动); RL 可在线但漂移衰减旧经验 &
仅 behavioral (rollout 间接比较) &
SWE-Gym [Pan24]; \S5.2 \\
\midrule
%── P6 ───────────────────────────────────────────────────
P6 Evaluator-Driven Optimization &
gradient training (RLHF/DPO) &
高·叠加 P4 成本 &
\textbf{保留}: 偏好迁移入 policy /
\textbf{注入}: trajectory-only 学不到的细粒度偏好 /
\textbf{损失}: 经 P4+P6 双重压缩 &
online/iterative·漂移敏感 &
behavioral·循环验证风险
(Evaluator $\rightleftarrows$ Policy) &
Constitutional AI [Bai22b]; \S6.1 \\
\midrule
%── P7 ───────────────────────────────────────────────────
P7 Parametric Externalization &
从 policy 采样生成 &
中·依赖采样质量·须先有 trained policy &
\textbf{保留}: 权重中的能力经采样部分恢复为 token /
\textbf{损失}: 采样未覆盖部分 /
\textbf{注入}: hallucination 风险 &
可增量生成新样本·跟随源模型更新 &
产物可读但不可追溯权重成分·完备性不可验证 &
Explorer [Pah25]; \S6.2 \\

\end{longtable}

} % end \small

\end{landscape}

\textbf{Production Mechanism} 把七条路径沿是否更新权重切成两组。P1、P2、P3 的 cache 子类与 P7 不更新权重，其中 P1、P2 靠 LLM inference，P3-cache 靠 gradient-free 的前向状态留存，P7 靠从已训练模型采样；P3 的 trained 子类与 P4、P5、P6 经 gradient training，需数据筹备与优化的整套机制，单步资源消耗高于一次 inference。P2 在 inference 之外可叠加 execution + verification loop，是不更新权重一侧唯一带客观产出闸门的机制。

\textbf{Construction Cost} 在 gradient-based training 一类整体高、LLM inference 一类整体低。P4 卡在标注或明确成败信号的获取，P5 卡在大规模高质量 trajectory，P6 在 P4 的 Evaluator 训练之上再叠一层、成本累加。P1 可全自动、低成本，P3-cache 获取成本极低，P2 中等、需可执行环境；P7 受制于采样质量，且须先有 trained policy——它的低成本只就外化这一步而言，不含上游训练。

\textbf{Information Transformation} 的关键区分在 loss 的两种性质。七条路径注入的新信息都来自转化模型的先验、而非源经验。P1（全局的细节信息）、P2（丢 NL nuance）、P3（丢符号可读性）的损失均为 lossy-by-design，损失类别可控、由抽象策略调节；P4、P5 在 by-design 之外混入 lossy-by-limitation——P4 的标注噪声与标注者不一致直接进权重，P5 的 SFT 欠/过拟合与 RL reward hacking 不可控，属工程风险而非设计意图；P6 的损失是 P4 阶段与 P6 阶段的双重压缩，原始经验语义到达后已不可辨识；P7 的损失来自采样未覆盖的部分，并注入 hallucination 风险。关于 traceability：P3-cache 的 KV 直接来自真实前向传播、能对应具体历史片段，P3-trained 的 soft prompt 已难以对应回单条经验，P4/P5/P6 把单条贡献经 batch 梯度平均进权重、无法追问某个决策用了训练集中哪条轨迹；P7 产物虽是 Tokenized 端的离散 token，却采样自 Parametric 权重、追不到源知识。这是 batch 梯度平均的结构性后果，与工程精度无关。

\textbf{Incrementality} 在 inference 一类天然增量——P1、P2、P3-cache 产生的经验可支持增加，但随着规模的变大，有效利用难度上升。P7 随源模型更新持续生成新样本；gradient 一类须 batch 重训，P3-trained 改训练数据即重训整个 bank、P4 的新标注需重训以保校准。P5 可热启动、RL 还能在线更新，但旧经验的效用随分布漂移衰减，在线性以遗忘旧经验为代价；P6 的 online/iterative 更新同样对漂移敏感。P1 一侧的 dynamic store 在 append 之外带 merge/prune 级联，单条写入可能触发已有条目的连带改写。

\textbf{Output Verifiability} 取决于产物容许哪种独立检验。P2 的 code skill 由环境 execution 给出客观二值信号（如 Voyager 的 verification gate，通过/失败可判定），P1 的经验条目和P7的生成数据只能靠人或强 LLM 主观判断。Parametric 端，P4 的 Evaluator 尚可用 hold-out calibration 间接核验，P3、P5、P6 退到只能借 rollout 或 ablation 推断。

P6 与 P7 都以已固化的 Parametric 经验为源、做第二步加工，方向相反、在 pathway 图中互补：P6 把 Evaluator 的判断进一步写入 Policy 权重、深入 Parametric 端，P7 把 Policy 或 Evaluator 的能力采样回 Tokenized、退出 Parametric 端，使权重中的隐式经验重新成为可供其他路径使用的显式载体。循环验证与 hallucination 是仅现于转化过程的两类风险：前者来自 P6 中 Evaluator 与 Policy 互为参照、缺独立外部 ground truth（ECHO \cite{Li26l} 的同步更新可缓解、不能消除），后者来自 P7 从权重采样、缺环境 grounding 时产出似真而无据的 trajectory。

% Policy 由 P5、P6 两条路径产生，二者代价画像不同，使 Section~\ref{sec:selection_constraints} 对 Policy 目标的可行性判断成为 P5 还是 P6 的选择、而非二元判断；

\subsection{Carrier Utilization Profiles}
\label{sec:util_profiles}

利用画像刻画五类载体生产完成后如何进入决策循环。载体的进入方式由其存在形态决定，与生产它的路径无关：Latent 向量无论来自 cache 留存还是训练所得模块，都作为连续向量参与 forward pass。沿四个维度展开。Access Interface（D1）：经验从 $c_t$ 还是 $\theta$ 进入决策循环。Access Granularity（D2）：能否在推理时有选择地激活特定经验，还是整个载体始终全开。Utilization Cost（D3）：使用经验需支付什么资源、计算发生在何处。Revisability（D4）：经验有误时如何修正，修正是否波及其他经验。

% Schematic artifact 无论是 code、workflow 还是 graph，都是 $c_t$ 中的强形式化离散 token、靠结构被操作（parse、execution 或 traversal），其中 code-skill 形态把执行卸至外部 executor。

% 这些层内差异改变具体单元格的取值（利用代价、修正方式），但不改变载体的可寻址性与访问粒度，记为单元格内限定，不另立载体。

% 可解释性与控制自由度都是访问粒度的下游读数。可读性随粒度与层次走——Tokenized 因逐条可寻址加离散符号而可读，Parametric 因不可寻址加连续权重而不可读；控制自由度同样系于粒度——逐条可寻址允许在检索、编排、注入各环节调节，不可寻址则不允许。二者作为任务约束的对偶留待 §8.3。

% 所需宏包：
% \usepackage{booktabs}
% \usepackage{tabularx}
% \usepackage{ctex}        % 中文支持（XeLaTeX / LuaLaTeX）

\begin{table}[htbp]
  \centering
  \caption{Carrier utilization profiles across four dimensions.}
  \label{tab:8.2}
  \small
  \renewcommand{\arraystretch}{1.5}
  \begin{tabularx}{\textwidth}{lXXXX}
    \toprule
    \textbf{Carrier} &
    \textbf{Access Interface} &
    \textbf{Access Granularity} &
    \textbf{Utilization Cost} &
    \textbf{Revisability} \\
    \midrule
    Narrative &
    context injection（$c_t$） &
    per-item &
    Retrieval Cost + Inference Cost &
    - \\
    \midrule
    Schematic &
    context injection（$c_t$） &
    per-item &
    Executor Cost + Inference Cost &
    - \\
    \midrule
    Latent &
    forward-pass attention &
    per-item &
    Inference Cost &
    - \\
    \midrule
    Policy &
    forward-pass weights（$\theta$）$\rightarrow$ $a_t$ &
    always-on &
    Inference Cost &
    - \\
    \midrule
    Evaluator &
    forward-pass weights（$\theta$）$\rightarrow$ 候选信号 &
    always-on &
    Inference Cost &
    - \\
    \bottomrule
  \end{tabularx}
\end{table}


访问接口由存在形态锁定：Tokenized 走 context injection（$c_t$），Latent 走 forward-pass attention（仍在输入侧，但以连续向量参与），Parametric 走 forward-pass weights（$\theta$）。选定载体即锁定接口。访问粒度沿 Tokenized → Latent → Parametric 单调变粗：Tokenized 逐条取用（如取单条 reflection 或 skill），Latent 介于逐条和整体之间，Parametric 全局生效。

层内的两个分叉不改变这条单调性，各自只动一两个维度。Tokenized 内，Narrative 与 Schematic 同处 per-item、同样可读，差别仅在复用机制——Narrative 靠 language understanding 阅读，Schematic 靠 parse、execution 或 graph traversal 按结构操作；这一差别落在 Output Verifiability（§8.1）与利用延迟（executable Schematic 的执行同步阻塞），不动可寻址性。Parametric 内，Policy 与 Evaluator 共享同一份存储画像——都不可寻址、都 retrain 才能改、复用都零边际——差别纯在功能角色：Policy 的 forward pass 直接产 $a_t$，Evaluator 的产一个关于候选的信号，可被训练消费（→ Evaluator-Driven Optimization）或推理消费（对 $\pi_\theta$ 候选 rerank、Best-of-N 筛选、tree search 剪枝）。以 PRM \cite{Lig23} 为例，同一套 step-level score 既能作 RL 的 dense reward，又能作 Best-of-N 的 step 级筛选或 tree search 的剪枝准则——这种下游灵活是"评估型输出可被多处取用"的性质，选哪个 evaluator、信号注入到哪个环节、取 outcome 还是 step 粒度，是把功能模块接入决策环的编排（属 §7），不是 Evaluator 的存储比 Policy 更可寻址。

在复用代价这一维上，Tokenized 与 Parametric 两端的取值难以：Tokenized 的经验可以单条取用，但是每次需要重新载入 Context 进行计算，存在复用成本，且当经验库的规模较大时，存在检索成本。Latent 居中，以 Hidden States 的形式参与前向计算，但同时也存在取用成本，比如检索或生成特定的 Hidden States。Parametric 的经验全局生效，每次复用边际成本可认为为零。

\subsection{Constraint-Driven Selection}
\label{sec:selection_constraints}

经验的载体与转化路径由任务的约束界定。将任务约束沿5个轴展开：Build Cost （C1） 刻画转化经验带来的成本；Serving Cost（C2）刻画使用经验的代价；Inspectability（C3）刻画经验的可解释性，包含 readability 和 verifiability；Maintainability（C4）刻画经验的维护成本，包含错误经验的修正 （error correction）和环境的非稳定性变化（non-stationarity） ；Scale（C5）刻画经验的可扩展性；


\begin{table}[htbp]
  \centering
  \caption{Constraint axes, their sub-constraints, and the \S8.1/\S8.2 dimensions each engages.}
  \label{tab:8.3}
  \begin{tabularx}{\textwidth}{llXX}
    \toprule
    约束轴 & 子约束 & 利用侧维度（\S8.2） & 生产侧维度（\S8.1） \\
    \midrule
    Build Cost
      & training resources
      & ---
      & Construction Cost、Mechanism \\
    Serving Cost
      & latency / context budget
      & Utilization Cost、Access Interface
      & --- \\
    Inspectability
      & readability
      & Access Granularity
      & Information Transformation（traceability） \\
    Inspectability
      & verifiability
      & ---
      & Output Verifiability \\
    Maintainability
      & error correction / non-stationarity
      & Revisability
      & Incrementality \\
    Scale（放大轴）
      & volume / arrival rate
      & 放大各轴
      & 放大各轴 \\
    \bottomrule
  \end{tabularx}
\end{table}


\begin{table}[htbp]
  \centering
  \caption{Directional pressure from constraint sub-dimensions to carrier classes. $\mathtt{++}$\,强烈偏向，$\mathtt{+}$\,偏向，$\circ$\,中性，$\mathtt{-}$\,不利，$\mathtt{{-}{-}}$\,强烈不利,指载体在该约束下的适配度。}
  \label{tab:8.4}
  \begin{tabular}{ll*{5}{c}}
    \toprule
    约束轴 & 子约束 & Narrative & Schematic & Latent & Policy & Evaluator \\
    \midrule
    \multirow{2}{*}{Serving Cost}
      & latency
        & $-$ & $--$ & $\circ$ & $++$ & $++$ \\
    \cmidrule(l){2-7}
      & context budget
        & $-$ & $--$ & $+$ & $++$ & $++$ \\
    \midrule
    \multirow{2}{*}{Inspectability}
      & readability
        & $++$ & $++$ & $-$ & $--$ & $--$ \\
    \cmidrule(l){2-7}
      & verifiability
        & $++$ & $++$ & $-$ & $--$ & $--$ \\
    \midrule
    \multirow{2}{*}{Maintainability}
      & error correction
        & $++$ & $+$ & $-$ & $--$ & $--$ \\
    \cmidrule(l){2-7}
      & non-stationarity
        & $++$ & $+$ & $-$ & $--$ & $--$ \\
    \midrule
    Build Cost
      & training resources
        & $+$ & $++$ & $-$ & $--$ & $--$ \\
    \bottomrule
  \end{tabular}
\end{table}

\textbf{Build Cost} 由转化是否更新权重决定，资源含算力与标注两类。gradient-free 的转化——P1、P2 与 P3 的 cache 子类——预算受限时仍可执行，不付训练成本；gradient-based 的转化——P3 的 trained 子类与 P4、P5、P6——需完整训练循环，预算不足时不可行，Latent 在此沿 cache 与 trained 分化最明显。算力与标注可分别稀缺：纯环境可验证信号（unit test、ground-truth match）驱动的 Policy RL 不需人工标注，而依赖人类示范或偏好标签的 SFT 与 reward modeling 在标注稀缺时受限更重。

\textbf{Serving Cost} 分时间与空间两面，区别在该代价随每次复用重新产生、还是一次性摊销。Narrative 两面均付、且每次复用重付：条目须检索并 prefill 进 context \cite{Shi23b, Zha23c}。Schematic 分形态——executable 将计算卸至外部 executor、不占 context，但执行同步阻塞决策；non-executable 仍注入 context。Latent 较为特殊，最终的复用方式都是以向量的形态直接参与模型的前向计算，但是其复用成本即可通过检索已有记忆库选出，也可通过相关模块来临时生成。Policy 与 Evaluator 边际复用成本为零，权重常驻。

\textbf{Inspectability} 分为 readability 与 verifiability，前者要逐条可读，后者要使用前可验证。Narrative 逐条可读、可主观核验，两者俱强 \cite{Zha23c}。Schematic 分形态：non-executable 的描述性结构可读性高，可主管验证 \cite{Ano24}；executable 的程序需技术能力方能阅读、可读性相对较低，但其执行可在落库前客观测试 \cite{Wan23c}。Latent， Policy 与 Evaluator 既不可读也无法事前验证。

\textbf{Maintainability} error correction 是修正已知错误的速度，non-stationarity 是跟随漂移的能力。Tokenized 一般逐条独立，单条修正不触发全局重写，新增条目即时可检索 \cite{Zha23c,Wan23c}。 对错误的修正较快，且对环境变化的适应能力较好。Latent 两面俱弱于 Tokenized。 Parametric 的表现最差，修正需 retrain 并可能引发 catastrophic forgetting，其在线更新以近端经验覆盖旧经验、旧经验效用随之衰减 \cite{Bai24}。


\textbf{Scale} 经验体量放大其余各轴的压力。随着经验库的增大，Tokenized 的检索质量会下降 \cite{Fu24}，维护成本会非线性的增加，加入 Context 后的推理成本接近是固定的；Parametric 推理成本与训练数据体量解耦、训练完成后保持恒定 \cite{Pan24}，其构建成本与维护成本会线性增加。 Latent 既存在检索使用的方式，有存在需要由训练好的模型临时生成的方式，介于两者之间。


约束为对应维度设定阈值，载体能否进入决策循环，取决于其取值是否满足。约束相容时逐一剪枝即可定位可行载体；约束冲突时单载体无解，迫使载体在时间或空间上组合。约束还会被误判、随部署漂移，使最优配置成为时间的函数。
\section{Open Problems and Future Directions}
\label{sec:open_problems}



%%
%% The acknowledgments section is defined using the "acks" environment
%% (and NOT an unnumbered section). This ensures the proper
%% identification of the section in the article metadata, and the
%% consistent spelling of the heading.
\begin{acks}
To Robert, for the bagels and explaining CMYK and color spaces.
\end{acks}

%%
%% The next two lines define the bibliography style to be used, and
%% the bibliography file.
\bibliographystyle{ACM-Reference-Format}
\bibliography{sample-base}


%%
%% If your work has an appendix, this is the place to put it.
\appendix

% \input{sections/appendix}

\end{document}
\endinput
%%
%% End of file `sample-manuscript.tex'.
